\documentclass[11pt]{article}
\usepackage[margin=0.9in,top=0.85in,bottom=0.85in]{geometry}
\usepackage{amsmath,amssymb}
\usepackage{graphicx}
\usepackage{float}
\usepackage{xcolor}
\usepackage{framed}
\usepackage{enumitem}
\usepackage[hidelinks]{hyperref}
\usepackage{parskip}

\definecolor{shadecolor}{RGB}{239,236,227}
\definecolor{sectblue}{RGB}{27,72,120}
\definecolor{stucol}{RGB}{27,90,140}
\definecolor{klecol}{RGB}{150,40,30}
\definecolor{beatgray}{RGB}{110,110,110}

\setlength{\parindent}{0pt}
\setlength{\parskip}{4pt plus 1pt}

% ---- dialogue ----
\newcommand{\Stu}[1]{\par\smallskip\noindent\hangindent=1.9em\hangafter=1%
  \textcolor{stucol}{\textbf{Student.}}~#1\par}
\newcommand{\Kle}[1]{\par\smallskip\noindent\hangindent=1.9em\hangafter=1%
  \textcolor{klecol}{\textbf{Klein.}}~#1\par}
\newcommand{\walkbeat}[1]{\par\smallskip{\centering\small\itshape\textcolor{beatgray}{#1}\par}\smallskip}

% ---- figures ----
\newcommand{\walkfig}[2]{%
  \begin{figure}[H]\centering
  \setlength{\fboxsep}{0pt}\setlength{\fboxrule}{0.4pt}%
  \fcolorbox{black!25}{white}{\includegraphics[width=0.97\linewidth,height=0.56\textheight,keepaspectratio]{img/#1}}\\[4pt]
  {\footnotesize\itshape #2}
  \end{figure}}

% ---- Klein's notebook ----
\newenvironment{knote}%
  {\begin{snugshade}\small\noindent\textbf{\textcolor{klecol}{Klein's notebook.}}\itshape\ }%
  {\end{snugshade}}

\renewcommand{\thesection}{\Roman{section}}
\begin{document}
\thispagestyle{empty}
\begin{center}
{\LARGE\bfseries Knots \& Primes --- a walkthrough}\\[6pt]
{\large A graduate student demonstrates the app to Professor John Klein}\\[14pt]
{\normalsize Full transcript, tab by tab, with Klein's questions, criticisms, and notebook}\\[4pt]
{\footnotesize Seventy-six screenshots of the live application. Every number visible in them was
computed in the browser at page load or on the click shown; nothing is typeset by hand.}\\[10pt]
{\footnotesize\itshape \today}
\end{center}
\vspace{10pt}
\begin{snugshade}\small
\textbf{How to read this document.} The student drives; the professor interrupts. Each figure is a
region of the running app, captured in the state the student had put it in --- presets clicked, cards
opened, the certificate cockpit run on Klein's own triple, the self-test suite executed live. Shaded
boxes are Klein's notebook, written as the demonstration proceeds. The closing post-mortem is the
debrief afterwards, and ends with a prioritized punch list of what needs to be improved.
\end{snugshade}
\vspace{6pt}
\hrule
\vspace{12pt}

\section*{The walkthrough}
\walkbeat{Klein sits, pulls the laptop a quarter turn toward himself, and does not take his coat off.}

\Kle{Before you show me a single picture. What is this, in one sentence, and what is it not?}

\Stu{It is one HTML file that runs Mazur's knots-and-primes dictionary as a working exhibit: a single story --- an intersection count, what survives when the count vanishes, and the zeta functions that record the whole hierarchy --- told simultaneously in topology and in number theory, at the smallest examples where each level is genuinely nontrivial. The Hopf link and the pair $(5,13)$; the unlink and $(5,29)$; the Borromean rings and the triple $(13,61,937)$. What it is not: a theorem, and not a functor. Nothing here sends links to sets of primes.}

\Kle{Good, because the functor question was going to be my first one. House rules.}

\Stu{Three, and they are enforced rather than promised. Every number on the screen is computed at page load or on a click --- there are no typed-in constants. Everything structural is cited by name and year at the point of use. And the three words \emph{theorem}, \emph{computed live}, and \emph{analogy} never trade places; where the app can only manage the third, it says so in the same sentence, not in a footnote. There is a button in the footer that recomputes every number and also asserts the rendered prose against the live page --- one hundred fourteen checks, sixty-two computational and fifty-two text pins.}

\Kle{Fifty-two text pins means you expect your own prose to rot faster than your arithmetic. That is more honest than most of what I read. Go.}

\subsection*{Dictionary}

\walkfig{d1-hero}{The header and the eight-tab bar: Dictionary, Pairs, Triples, Ladder, Quadruples, Zeta, Q and A, Experimental.}

\Stu{Eight tabs, and the order is the argument. Dictionary states the correspondence. Pairs computes the intersection count at the two-component level. Triples computes the first invariant that only exists because the pairwise count vanished. Quadruples goes one further. Zeta cashes all of it out. Q and A is fourteen objections with answers, and Experimental is the frontier, walled off --- nothing on it is load-bearing for the other seven.}

\Kle{Then why is Ladder sitting between Triples and Quadruples instead of at the end with the other machinery?}

\Stu{Because it is used there, not proved there. The Ladder tab is the unipotent tower over $\mathbb{F}_2$ --- $N_2$, $N_3$, $N_4$, of orders $2$, $8$, $64$. The superdiagonal carries the pairwise level; each deeper diagonal is only defined once everything above it vanishes. Triples needs $N_3(\mathbb{F}_2)$, which is $D_4$; Quadruples needs $N_4(\mathbb{F}_2)$. Putting the ladder after both would mean two tabs quietly borrowing a group the reader has not met.}

\Kle{Fine. That is a defensible reason and you gave it in one breath, which is more than I expected. The subtitle says dotted terms carry definitions. So the page has two layers.}

\Stu{Yes, and that is a real tension I will not pretend away. Hover, Tab, or click pins them; Escape closes.}

\walkfig{d2-whyknots}{The ``why is a prime a knot'' box, opened; Artin-Verdier duality sits under a dotted rule.}

\Stu{This is the load-bearing paragraph of the whole app. Each prime is secretly a circle: $\mathrm{Gal}(\overline{\mathbb{F}_p}/\mathbb{F}_p)$ is $\hat{\mathbb{Z}}$, which is exactly the profinite fundamental group of $S^1$, so $\mathrm{Spec}\,\mathbb{F}_p$ is a circle in disguise. And the ambient space behaves three-dimensionally: etale cohomological dimension $3$. A circle in a $3$-space is a knot. That is Mazur's observation, and it is the reason the table below is a dictionary rather than a coincidence.}

\Kle{Stop. Cohomological dimension three is a number, not an orientation. Poincare duality in top degree needs a fundamental class and it needs the thing to be closed. $\mathrm{Spec}\,\mathbb{Z}$ is affine and it is not compact. Where is your compactification, and what is your orientation sheaf?}

\Stu{Both are in the tooltip on Artin-Verdier, and I will pin it rather than paraphrase it. The trace isomorphism is $H^3(\mathrm{Spec}\,\mathcal{O}_K, \mathbb{G}_m) \cong \mathbb{Q}/\mathbb{Z}$, and the pairing of $H^i(F)$ against $\mathrm{Ext}^{3-i}(F, \mathbb{G}_m)$ into $\mathbb{Q}/\mathbb{Z}$ is perfect for constructible $F$. So $\mathbb{G}_m$ is the orientation sheaf and that degree $3$ is the precise sense in which $\mathrm{Spec}\,\mathbb{Z}$ behaves like a closed oriented $3$-manifold. The fine print is in the same tooltip, in the same size type: the real place introduces $2$-torsion corrections, and they are handled by compactifying $\mathrm{Spec}\,\mathbb{Z}$. Cited, not computed --- Artin-Verdier, Woods Hole 1964; Mazur, Ann. Sci. ENS 6 (1973).}

\Kle{Then that tooltip is not a tooltip. It is the justification for your entire premise, and until I put a cursor on it, the page asserts ``behaves 3-dimensionally'' with a parenthesis. Anyone who reads this without hovering has read an unsupported claim in a serif font. Promote it.}

\Stu{That is fair and I will not argue it. The visible sentence is doing rhetorical work the hidden sentence is doing mathematical work.}

\begin{knote}
Premise: $\mathrm{Spec}\,\mathbb{F}_p \simeq S^1$ via $\hat{\mathbb{Z}} = \pi_1^{et}(S^1)$; ambient etale cd $= 3$; A-V duality with $\mathbb{G}_m$ as orientation sheaf, $H^3 \cong \mathbb{Q}/\mathbb{Z}$. Real place gives 2-torsion, killed by compactifying. He knows the fine print and it is written down --- but it is written down where I have to go looking. Content: correct. Typography: it hides its own load-bearing wall.
\end{knote}

\walkfig{d3-table}{The full dictionary table --- topology column against arithmetic column, eleven rows.}

\Stu{Eleven rows, read across. Link in $S^3$ against a set of primes in $\mathrm{Spec}\,\mathbb{Z}$, each $\equiv 1 \pmod 4$. Seifert surface against quadratic resolvent $k_i = \mathbb{Q}(\sqrt{p_i})$, ramified only at $p_i$ --- touching only its own knot. Puncture point against non-split prime, that is Frobenius. Linking number against Legendre symbol. Symmetry of lk against quadratic reciprocity. Then the vanishing rows: surface intersection $\gamma = \Sigma_1 \cap \Sigma_2$, closed since lk $= 0$, against the Redei element $\alpha$, which exists because $(p_1/p_2) = +1$. Triple linking $\bar\mu(123)$ against the Redei symbol. Then Alexander polynomial against Iwasawa polynomial, torsion against Dedekind zeta, orbits against primes, eigenvalues against characters.}

\Kle{The $\equiv 1 \pmod 4$ in row one. You are telling me the arithmetic side gets to choose its objects and the topology side does not.}

\Stu{Yes, and the note under the table says exactly what it buys: the congruence tames ramification at $2$ and at $\infty$ in the quadratic resolvents, and it plays the role of orientability. The same note says the taming does not propagate up the tower.}

\Kle{Now the row I actually want to fight about. Seifert surface against quadratic resolvent. A knot bounds infinitely many Seifert surfaces and no canonical one. $\mathbb{Q}(\sqrt{p})$ is unique on the nose. Those are not the same kind of object and you have put them in the same row.}

\Stu{The note under the table concedes precisely that, in those terms: a knot bounds many Seifert surfaces and every construction has to be checked independent of the choice, whereas the resolvent is unique, so on the arithmetic side well-definedness is free. It is listed as an asymmetry, not smoothed over.}

\Kle{Good. That paragraph is the best writing on the tab, and it is set in the smallest type on the tab. Second complaint, and this one is pure carelessness --- your first row reads $K_1 \cup \cdots \cup K_r$ with a typeset subscript on the one and a literal underscore on the $r$. Same in the arithmetic cell, $p_1$ typeset, $p_r$ with an underscore. It looks like the file leaked its own source.}

\Stu{It did. Those two cells are plain text where the neighbours are marked up. That is a bug, not a decision.}

\begin{knote}
Table is honest where it costs it something: the Seifert/resolvent non-uniqueness asymmetry is stated, not buried --- credit. But $K_-r$ and $p_-r$ render as raw ASCII subscripts beside properly set ones. Whole tab loses authority on a typo. Tell him to fix it before anyone else sees it.
\end{knote}

\walkfig{d4-thread}{The $\zeta \leftrightarrow \Delta$ thread box: ``Where this is going.''}

\Stu{This box is a running device. Any box wearing the $\zeta \leftrightarrow \Delta$ badge carries the same thread across tabs. The claim is: two zetas --- $\zeta_K$ of a number field against the Lefschetz zeta of a monodromy --- and two polynomials, $\Delta_L$ against $f_S$, each polynomial being the zeta of its own cover up to units. Milnor-Turaev on the left, the Main Conjecture on the right. And two dials feed those zetas: linking, which is Taylor coefficients at $t = 1$ and is this app's whole story, and knotting, which is eigenvalues pushed off $1$ and is not.}

\Kle{Separating those two dials is the right instinct. Most expositions of this material blur them and then wonder why the analogy stops paying.}

\Stu{The Zeta tab has contrast cards that keep them apart deliberately --- the figure-eight knot against $\mathbb{Q}(\sqrt{5})$, where $\varphi^2$ shows up as both a stretch factor and a fundamental unit, and the app labels that coincidence as a coincidence.}

\Kle{Then here is your problem in this box. Milnor-Turaev gets two names. ``The Main Conjecture'' gets no name, no year, and no article. On a page whose stated rule is name-and-year at the point of use, that is a violated rule, and it is the one row where I have the least prior --- I do not do Iwasawa theory.}

\Stu{You are right that it is violated here. The names exist in the app: Mazur-Wiles, Invent. Math. 76 (1984), and $\ell = 2$ --- which is the case every computation in this app instantiates --- is Greither 1992. Both are in the tooltip on the Zeta tab and again in the Q and A tab. But they are two clicks from the sentence that needs them.}

\Kle{Two clicks is not a citation. And note what you just told me: the only case your app ever actually computes is the case that was proved eight years later by somebody you did not name.}

\Stu{Conceded. That should be inline.}

\walkfig{d5-whyfs}{The $f_S$ why-box: finite deck group over $\mathbb{Q}$, unramified outside $S \cup \{\infty\}$, $\mathbb{Z}/4 \times \mathbb{Z}/4$, and the totally real $\mathbb{Z}/4 \times \mathbb{Z}/2$.}

\Stu{This box exists because you would have asked. On the knot side, $\Delta_L$ lives over Laurent polynomials with deck group $\mathbb{Z}^r$. Over $\mathbb{Q}$ there is no rank-$r$ tower mirroring that. For $\ell \notin S$, the Galois group of the maximal abelian $\ell$-extension of $\mathbb{Q}$ unramified outside $S \cup \{\infty\}$ is a product of $\mathbb{Z}/\ell^{e_i}$ with $e_i = v_\ell(p_i - 1)$, and that is \emph{finite}. For $S = \{13, 61\}$ and $\ell = 2$ it is $\mathbb{Z}/4 \times \mathbb{Z}/4$.}

\Kle{And the $\cup \{\infty\}$ is doing what, exactly? Because in my world you do not get to append a point to your ramification set and keep a straight face.}

\Stu{It is a convention, the box says so, and it says where it bites: at $\ell = 2$. Complex conjugation lands on the order-two element $(2,2)$, so the totally real tower --- unramified at $\infty$ --- is the strictly smaller $\mathbb{Z}/4 \times \mathbb{Z}/2$. And it spells out why: $p \equiv 1 \pmod 4$ tames $\infty$ for the resolvent $\mathbb{Q}(\sqrt{p})$ only, not higher up. The quartic subfield of $\mathbb{Q}(\zeta_{13})$ is already imaginary.}

\Kle{So the orientability convention you sold me in row one of the table does not survive one floor up the tower.}

\Stu{Correct, and the note under the table says that too --- ``the taming does not propagate up the tower.'' The consequence is the honest one: $f_S$ lives over the finite deck group's ring, while Iwasawa's own $\Lambda = \mathbb{Z}_\ell[[T]]$ series runs over the cyclotomic tower, whose deck group really is $\mathbb{Z}_\ell$ but which ramifies at $\ell$ itself. Two different rings. The Q and A tab lists that as fracture four of five in the dictionary.}

\Kle{Then the $\Delta_L$ against $f_S$ row of your table is the weakest row in the table, and you know it, and you say so somewhere else. Put a marker in the row.}

\Stu{Yes.}

\Kle{Last thing on this box, and it is the same complaint as before but worse. Look at it. $\mathbb{Z}_\ell[[u_1,\ldots,u_r]]$ with underscores and square brackets, $v_\ell(p_i - 1)$ with an underscore, and $(1+u_i)$ raised to $\ell^{e_i}$ written with a caret. This paragraph is the most technical thing on the tab and it is typeset like a mailing-list post while the table above it is set properly.}

\Stu{Same defect as the table row, larger surface. The box was written last and it shows.}

\begin{knote}
$f_S$ box: the real content of the tab. Deck group over $\mathbb{Q}$ is finite --- $\prod \mathbb{Z}/\ell^{e_i}$, $e_i = v_\ell(p_i-1)$; $S=\{13,61\}$, $\ell=2$ gives $\mathbb{Z}/4 \times \mathbb{Z}/4$, totally real part $\mathbb{Z}/4 \times \mathbb{Z}/2$ because conjugation hits $(2,2)$. So the $\Delta_L / f_S$ row is a row between objects over different rings, and he says so --- but says it here and in Q and A, not in the row. Also: ASCII math leaking through the serif. Twice now.
\end{knote}

\subsection*{Pairs}

\walkfig{p1-lede}{The Pairs tab lede.}

\Stu{One intersection count, computed in both languages: lk by counting punctures through a spanning surface, $(p/q)$ by Euler's criterion. Then one symmetry with the same proof twice. And the link at the end of the lede goes straight to the bottom of the tab, where the count becomes a picture --- the double cover drawn as a lift diagram.}

\Kle{You put the destination in the first sentence. I approve of that and I want to see whether the tab earns it. What is computed here and what is asserted?}

\Stu{Every Legendre symbol on this tab is computed live by Euler's criterion --- $p^{(q-1)/2} \bmod q$, BigInt fast exponentiation, one function that every symbol in the app goes through. The Hilbert symbol at $v = 2$ is computed by enumeration. The topology side is a drawing plus cited reasoning, and it is labelled as cited at the two places where it matters.}

\walkfig{p2-canvas}{The two-circle canvas in the unlinked state, with the state label lk $= 0$.}

\Stu{$K_1$ blue, $K_2$ green, arrowheads marking orientations. In this state $K_1$ is consistently on top at both crossings, so the diagram is the unlink and the label reads lk $= 0$. The ``Lift apart'' button animates them apart and prints ``Separated --- the circles were never entangled.''}

\Kle{Two crossings, both the same sign resolution, and you are calling it the unlink. That is a diagram of the unlink, not the split unlink. Do you distinguish them anywhere, because your Alexander statement further down had better care.}

\Stu{It does, and the app makes the distinction in a tooltip on the phrase ``split unlink'': two circles in disjoint balls, separable by a sphere, and splitness is strictly stronger than being unlinked crossing by crossing. The polynomial claim in the lower strip is the split one --- $\Delta \doteq 0$ is the classical splitness test.}

\Kle{Good, the drawn state is isotopic to split so you are safe, but you are safe by an inch. Now the sub-caption. Read me the sentence about crossing signs.}

\Stu{``Crossing signs are read in the screen's y-down, left-handed frame, and the app only ever displays diagram counts that survive that choice --- parities and zero totals, never a bare handed sign.''}

\Kle{So you have made a design rule out of refusing to commit to an orientation. In a browser, with a y-down canvas, I will accept it, because the alternative is a mirror-image bug that nobody catches for a year. But say it once and loudly rather than in the small type under one figure.}

\walkfig{p3-hopf-canvas}{The same canvas toggled to the Hopf link; the label reads lk $= \pm 1$.}

\Stu{Toggle, and the crossings alternate --- one under each way. The label becomes lk $= \pm 1$, Hopf link. That $\pm$ is the design rule you just read, cashed out: the app will tell you the linking number is odd and nonzero, and it will not tell you the handedness, because the handedness is an artifact of which way the canvas $y$-axis points.}

\Kle{And ``Lift apart'' in this state?}

\Stu{It bounces --- the circles separate partway and come back, and the note reads ``Bounced back --- a genuine crossing survives every isotopy.''}

\Kle{That is a cartoon, not an argument. An animation that returns to its start proves nothing about isotopy; it proves your easing function has a turning point at fifty-five percent.}

\Stu{Agreed, and I will not defend it as more than a gesture. The actual invariance argument on this tab is the surface-independence why-box, which is a proof, and the product formula, which is a proof. The bounce is a mnemonic.}

\Kle{Then let me have the why-box, since you raised it.}

\Stu{``Two choices of Seifert surface glue along $K_1$ into a closed surface, against which the disjoint curve $K_2$ has algebraic intersection number $0$ --- so both puncture counts agree.''}

\Kle{Wrong as stated. You glue $\Sigma$ to $-\Sigma'$. If you glue two coherently oriented surfaces along their common boundary you do not get a cycle, you get something with a boundary counted twice. The orientation reversal is the entire content of the argument and it is missing from the sentence.}

\Stu{That is a real gap in the wording. The conclusion is right and the reason as written is not sufficient.}

\begin{knote}
Canvas: unlink and Hopf, toggled, arrowheads for orientation. Label refuses handedness on purpose --- y-down frame, stated rule, only parities and zeros displayed. Defensible in a browser. ``Lift apart'' bounce is decoration. Surface-independence why-box drops the orientation reversal: should read $\Sigma \cup (-\Sigma')$. Fix the sentence; the mathematics is fine underneath it.
\end{knote}

\walkfig{p4-calculator}{The Legendre calculator with the $(5,29)$ preset: both symbols $+1$, the reciprocity check, and the verdict banner.}

\Stu{Preset $(5,29)$. Euler's criterion gives $(5/29) = +1$ and $(29/5) = +1$, both computed, and the third row is a live check that they agree --- reciprocity, not assumed but verified against the two computed values. The banner says $+1$, $29$ splits in $\mathbb{Q}(\sqrt5)$: unlinked. Pressing the other preset gives $(5,13)$, both symbols $-1$, and the banner reads: the Hopf pair.}

\Kle{The reciprocity row. If both entries come out of the same function, checking that they agree is checking that your exponentiation is deterministic. It is not checking reciprocity.}

\Stu{They come out of the same function but from genuinely different computations --- $5^{14} \bmod 29$ and $29^{2} \bmod 5$. Nothing in Euler's criterion knows about the other direction, so the agreement is not automatic; a bug in the modular exponentiation would break both, but reciprocity is not smuggled in. The real independent check is one panel over, in the Hilbert table, and that one is built so that it can fail.}

\Kle{And what does pressing Compute do to the drawing on the left?}

\Stu{It syncs it, when the gloss is licensed --- if either prime is $\equiv 1 \pmod 4$ the app draws the unlink for $+1$ and the Hopf link for $-1$. Hold that thought, because in a minute I am going to show you the case where it deliberately does not sync, and you are going to dislike what that looks like.}

\walkfig{p5-hopf-thread}{The Pairs $\zeta \leftrightarrow \Delta$ thread box for $(5,13)$: Frobenius on the roots of $x^2 - 5$, one orbit of period $2$, and the local factor.}

\Stu{This is the thread badge again, now carrying a live computation. Frobenius $x \mapsto x^{13}$ acts on the two roots of $x^2 - 5$. Since $5$ is not a square mod $13$ the roots live upstairs in $\mathbb{F}_{13^2}$, and Frobenius swaps them: $r^{13} = (r^2)^6 \cdot r = 5^6 \cdot r = -r$, because $5^6 \equiv 12 \equiv -1 \pmod{13}$, and that $12$ is computed, not quoted. One orbit of period $2$, so the local factor is $(1 - 13^{-2s})^{-1}$. With $t = 13^{-s}$, split gives $(1-t)^{-2}$ and inert gives $(1-t)^{-1}(1+t)^{-1} = (1-t^2)^{-1}$ --- these are literally the Euler factors of $\zeta_{\mathbb{Q}(\sqrt5)} = \zeta(s) L(s,\chi)$.}

\Kle{Now the second paragraph, the one in blue.}

\Stu{That is the topology side, and note how it opens: ``cited reasoning, not computed.'' In the double cover of the $K_1$-complement --- the cover this symbol mirrors --- the preimage of $K_2$ is two disjoint copies when lk is even, one connected curve double-covering $K_2$ when lk is odd. Two period-one orbits against one period-two orbit. Same bookkeeping.}

\Kle{Two period-one orbits of \emph{what}? You have a deck group of order two acting on a preimage. Say which action.}

\Stu{The deck action, and the box says so --- ``two period-1 orbits of the deck action.'' The tooltip on the phrase gives the argument: $K_2$ lifts to a loop exactly when its class dies under $\pi_1 \to \mathbb{Z}/2$, that is when lk$(K_1,K_2)$ is even, and the components of the preimage are cosets of the image subgroup.}

\Kle{Then here is the real question, and it is the one your whole tab turns on. lk lives in $\mathbb{Z}$. Your symbol lives in $\{\pm 1\}$. Your cover only ever sees lk mod $2$. So the Hopf link and a link with lk $= 3$ are indistinguishable to everything on this tab, and you have drawn a dictionary in which one side forgets almost everything.}

\Stu{That is exactly right, and the app states it as a limitation rather than defending it. There is a why-box on the arithmetic panel: ``a curve can wind any number of times around a surface, while a degree-$2$ extension admits only one nontrivial Frobenius. The Legendre symbol is a mod-$2$ linking number.'' It is not a lossless correspondence, and where the app displays a topological count derived from the symbol, it derives it from lk parity, not from the symbol's value.}

\Kle{Then say ``mod-2 linking number'' in the table row on the Dictionary tab, not in a collapsed box three tabs away. Right now the table shows me $\mathbb{Z}$ against $\{\pm 1\}$ and lets me notice the gap myself.}

\begin{knote}
Thread box, $(5,13)$: $5^6 \equiv 12 \equiv -1 \pmod{13}$ computed; inert; one period-2 orbit; factor $(1-13^{-2s})^{-1}$. Euler factors of $\zeta(s)L(s,\chi)$. Topology half explicitly flagged cited-not-computed --- he is disciplined about that line and I have now seen him hold it three times. The lossiness is real: symbol sees lk mod 2 only. He knows; the admission is one collapsed box away from the table that needs it.
\end{knote}

\walkfig{p6-arc}{The arc figure --- $\Sigma_1$ horizontal, $\Sigma_2$ vertical, the segment $\Sigma_1 \cap \Sigma_2$ between them --- with caption and note.}

\Stu{This is the topological half of the reciprocity proof. The caption opens by heading off the misreading: this is a three-dimensional picture, not a Venn diagram. $\Sigma_1$ is a flat disk in a horizontal plane, $\Sigma_2$ a flat disk in a vertical plane, threaded through each other like two chain links with soap films. Two transverse planes meet in a line, so the two disks meet in a segment --- and that segment is the arc. One end sits on $K_1$, where $K_1$ pierces $\Sigma_2$, counted by lk$(K_2,K_1)$. The other sits on $K_2$, where $K_2$ pierces $\Sigma_1$, counted by lk$(K_1,K_2)$. A compact oriented $1$-manifold has signed boundary count zero. One object, two ends, so the counts agree.}

\Kle{The argument is right and the staging sentence is genuinely useful --- I have watched three people misread that exact figure as a Venn diagram. But look at your own picture. The words ``$K_2$ pierces $\Sigma_1$ here'' are struck through by the blue ellipse. So is the ``$\Sigma_1 \cap \Sigma_2$'' label. The two labels that name the two ends of your global object are the two labels you cannot read.}

\Stu{They are drawn under the ellipse rather than over it. That is a paint-order defect and it is on the one figure that has to be legible.}

\Kle{Second thing, and this one is mathematical. Both your surfaces are flat disks. That is available to you only because both components are unknotted. If $K_1$ is a trefoil, $\Sigma_1$ has genus one and it is not flat and it is not a disk, and your picture has no way to say so.}

\Stu{Conceded --- the drawing hard-codes the smallest case. What I would say in the figure's defence is that the argument in the note does not use disks anywhere: it uses transversality and the fact that a compact oriented $1$-manifold has signed boundary zero, and both survive genus. But the reader only sees disks, so the reader is entitled to think disks are load-bearing.}

\Kle{Then put one sentence under it saying the picture is the genus-zero case and the argument is not.}

\walkfig{p7-hilbert}{The Hilbert product-formula table for $(5,13)$: $v = \infty$, $v = 2$ by mod-8 enumeration with the closed form as cross-check, $v = 5$, $v = 13$, and the product row.}

\Stu{The arithmetic half of the same proof. Five rows. At $v = \infty$ the symbol is $+1$ because both primes are positive. At $v = 2$ it is computed --- genuinely computed, by testing primitive solvability of $z^2 \equiv px^2 + qy^2 \pmod 8$ over all $512$ triples, discarding the all-even ones. Then $v \notin \{2,p,q,\infty\}$, then $v = 5$ giving $(13/5) = -1$, then $v = 13$ giving $(5/13) = -1$. Product $+1$, and the note says the formula collapses to $(5/13)(13/5) = 1$: reciprocity is the closedness of one global object.}

\Kle{Why mod eight and not mod sixteen?}

\Stu{Because odd squares are determined mod $8$ and lift $8$-adically by Hensel, which is the sentence in the ``why'' column. That justification is cited, not proved in the app.}

\Kle{And the closed form sitting next to it --- the one with $(-1)$ to the product of $(p-1)/2$ and $(q-1)/2$?}

\Stu{That is deliberately kept out of the vote. The app's own comment says why: that closed form is reciprocity-equivalent, so if it computed the entry, the table would be assuming what the table is proving. It is displayed only as an independent cross-check, and the check is built so it can print ``DISAGREES'' rather than always printing a tick.}

\Kle{That is the right architectural decision and I want it noted that you made it. Most people would have used the closed form and never noticed they had assumed the theorem. Now --- row three.}

\Stu{$v \notin \{2, p, q, \infty\}$, value $+1$, reason ``both are local units.''}

\Kle{That is one row standing for infinitely many places, and it is the only row on this table that is not computed. It is hard-coded. You have a table that presents itself as a computation and one of its five entries is a theorem you typed in.}

\Stu{Correct, and it cannot be otherwise --- there are infinitely many places, so no browser will enumerate them. What the row asserts is that for odd $\ell$ not dividing $pq$, both $p$ and $q$ are $\ell$-adic units and the symbol of two units at an odd place is trivial. That is standard and it is true, but you are right that the row is doing the same rhetorical trick I just complained about on the Dictionary tab: it sits in a computed table wearing computed clothes.}

\Kle{Then mark it. One word --- ``cited'' --- in that cell, and the table becomes honest instead of merely correct.}

\walkfig{p8-withheld}{The $(3,7)$ case: both primes $\equiv 3 \pmod 4$, the warn line, the flipped reciprocity check, and the thread box withholding the linking gloss.}

\Stu{Now the case I flagged. Enter $3$ and $7$ --- both $\equiv 3 \pmod 4$, outside the paper's convention. Watch what the app refuses to do. $(3/7) = -1$, $(7/3) = +1$, and the reciprocity row prints the \emph{flipped} identity with a tick: $(p/q) = -(q/p)$, both $\equiv 3 \bmod 4$, and it names the $-1$ as the $2$-adic factor. The banner says $-1$, $7$ is inert in $\mathbb{Q}(\sqrt3)$ --- and then, in bold red, ``No linking gloss here.'' Because with both primes $\equiv 3 \pmod 4$ the symbol is order-dependent: swap the inputs and the banner flips, while lk is symmetric. The Hopf-unlink dictionary is reserved for pairs with a prime $\equiv 1 \pmod 4$.}

\Kle{And the thread box at the bottom does the same thing.}

\Stu{It withholds the topology half entirely --- ``withheld for this pair'' --- while keeping the Euler factor. The Euler factor survives because it reads $7$'s splitting in the specific field $\mathbb{Q}(\sqrt3)$, fixed once and for all, where the ordering is part of the data rather than a symmetry that has to hold.}

\Kle{That distinction is sharp and it is the right one. Refusing to draw a conclusion is the hardest thing to build into a demo and you built it. Now look at the top left of your own screen.}

\walkbeat{He puts a finger on the canvas panel, which still reads lk $= \pm 1$ (Hopf link).}

\Kle{Your drawing still says Hopf link. It is sitting six inches from a banner that says the Hopf dictionary does not apply to this pair. Every reader who scans left to right gets the gloss you just spent a paragraph withholding.}

\Stu{That is a genuine defect and it is a consequence of the fix rather than an oversight. The Compute handler only re-syncs the topology panel when at least one prime is $\equiv 1 \pmod 4$ --- precisely so it does not draw a picture the banner is disowning. But not updating it leaves the previous pair's picture on screen, which is worse, because a stale correct-looking picture reads as an answer.}

\Kle{It should blank, or grey out, or print ``no topological reading for this pair.'' Silence and staleness look identical to a reader.}

\Stu{Yes. Refusing to answer and forgetting to answer render the same way, and only one of them is the design.}

\Kle{One more thing, and then I will let this case go. Your banner says $7$ is inert in $\mathbb{Q}(\sqrt3)$. What is the discriminant of $\mathbb{Q}(\sqrt3)$?}

\Stu{Twelve. So it ramifies at $2$ as well as at $3$, and the ``surface touching only its own knot'' reading fails.}

\Kle{Does the app say that?}

\Stu{In the mixed case it does --- if exactly one prime is $\equiv 3 \pmod 4$, the warn line spells out that $\mathbb{Q}(\sqrt p)$ has discriminant $4p$, is ramified at $2$ as well as $p$, and that only splitting and symmetry survive. In the both-$\equiv 3$ case, which is this one, the warn line talks about the sign flip and the $2$-adic factor and never mentions the discriminant. So the caveat is stated in the case where it is half-needed and omitted in the case where it applies to both fields.}

\Kle{Which is the case a reader is most likely to be confused by. Move the sentence.}

\begin{knote}
$(3,7)$: $(3/7) = -1$, $(7/3) = +1$, flipped reciprocity checked and the $-1$ named as the 2-adic factor. Banner and thread box both refuse the lk gloss --- order-dependent symbol against symmetric lk. That refusal is the best thing on the tab. Two defects at the same moment: (i) canvas keeps the stale Hopf picture because the sync is suppressed --- silence looks like an assertion; (ii) disc $= 12$ caveat appears in the mixed-parity warning and vanishes in the both-3 warning, which is backwards.
\end{knote}

\walkfig{p9-hilbert37}{The Hilbert table for $(3,7)$, with the $2$-adic factor equal to $-1$.}

\Stu{Same table, same code path, different congruence class. Now $v = 2$ computes to $-1$ --- from the same $512$-triple enumeration --- and the closed form independently gives $-1$ and agrees. $v = 3$ gives $(7/3) = +1$, $v = 7$ gives $(3/7) = -1$, and the product is still $+1$ with a tick. The note reads: the product formula collapses to $(p,q)_2 \cdot (3/7)(7/3) = 1$ with $(p,q)_2 = -1$, so $(3/7)(7/3) = -1$ --- the sign flip of this congruence class is itself part of the closedness of the same global object.}

\Kle{So the $-1$ is not an exception to your global principle, it is a term in it.}

\Stu{That is the whole point of computing $v = 2$ instead of assuming it. In the $\equiv 1 \pmod 4$ world the $2$-adic factor is $+1$ and drops out silently, and you can go a long time believing the product formula \emph{is} reciprocity. Run it at $(3,7)$ and the missing factor appears with a place attached to it. The convention on the Dictionary tab is not a simplification, it is a choice to work where one specific local term vanishes.}

\Kle{Good. That is the strongest paragraph you have said all session, and it is the one thing here I would actually use in a lecture --- the archimedean and dyadic bookkeeping made visible by choosing bad primes on purpose. Does the app ever say it that plainly?}

\Stu{Not in that many words. It says the pieces --- the warn line, the product note, the orientability remark under the dictionary table --- but a reader has to assemble them.}

\Kle{Then that is a paragraph you have not written yet, and it belongs at the top of this tab.}

\walkfig{p10-lift}{The double-cover lift diagram: two sheets, $2:1$ arrows, two lifted loops against one loop winding twice, with caption and the Galois-correspondence line.}

\Stu{The count becomes a picture. Centre: a two-box tower --- $X = S^3 \setminus K_1$ at the bottom, its double cover $\tilde X$ on top, with the edge labelled ``$2$ sheets, degree $2$: $\mathbb{Q}(\sqrt p)/\mathbb{Q}$.'' Left panel, lk even: the two sheets are drawn as two parallel lines and $K_2$ lifts to two separate loops, one on each sheet, projecting $2:1$ down to the single loop downstairs --- split, two primes above. Right panel, lk odd: one loop travels along the top sheet, passes down to the bottom sheet, and closes up --- one loop winding twice, inert, one prime of residue degree $2$. Either way loops times winding equals $2$, which is the order of the deck group.}

\Kle{Loops times winding equals the order of the group, and $efg$ equals the order of the group. You are drawing $e = 1$, $f \cdot g = 2$.}

\Stu{Exactly, and the app says it as $|G| = \mathrm{count} \times \mathrm{size}$ --- a loop with monodromy of order $d$ has preimage $|G|/d$ loops each winding $d$ times, and a prime with Frobenius of order $f$ has $g = |G|/f$ primes above it each of residue degree $f$, with $e = 1$ because we are unramified. One bookkeeping, two languages. That identity is what the Triples and Quadruples tabs scale up, to $8$ sheets and $64$.}

\Kle{The correspondence line at the bottom.}

\Stu{``The lattice of intermediate covers mirrors the subgroup lattice of the deck group --- the Galois correspondence, both sides --- here at its smallest: one $\mathbb{Z}/2$ on each side.'' And it is labelled cited, not computed, with Hatcher section 1.3 for covers and Morishita chapters 2 through 4 for the joint dictionary.}

\Kle{Which is the correct label, because at $|G| = 2$ there is nothing to compute --- the lattice is two boxes and a line. This figure is doing something the rest of the tab is not, though: it is the only place where I can see, in one glance, that the arithmetic statement and the topological statement are two readings of the same diagram. The tables tell me; this shows me.}

\Stu{That was the intent, and it is why the lede links straight to it.}

\Kle{Two complaints and I am done. First: your tilde. ``$\tilde X$'' renders with the accent floating half off the letter, in three separate labels, and on a figure whose entire job is to distinguish upstairs from downstairs, the upstairs symbol is the one that looks broken. Second, and this is the one that matters --- your left panel says ``lk even (split).'' Even, not zero. So you are correctly telling me the cover cannot distinguish lk $= 0$ from lk $= 2$.}

\Stu{Yes. The word ``even'' is chosen deliberately, and the count is derived from lk parity rather than from the symbol.}

\Kle{Then the picture is more honest than the table on the Dictionary tab, which puts $\mathbb{Z}$ opposite $\{\pm 1\}$ and lets it pass. Make the table as honest as the figure.}

\begin{knote}
Lift diagram: the one thing here I would put on a slide. Two sheets, $K_2$ lifts to 2 loops when lk even / 1 loop winding twice when lk odd; loops $\times$ winding $= 2 = |{\rm deck}|$, which is $efg$ with $e=1$. Same bookkeeping, two languages, one picture. Scales to 8 and 64 sheets on later tabs.

Running tally, two tabs in: mathematics checks everywhere I have pushed. Failures are all of the same species --- support exists but is one hover, one click, or one tab away from the sentence that needs it, and the page presents cited things and computed things in identical type. Also: raw ASCII subscripts twice, labels painted under a figure, and a stale canvas that reads as an assertion. He concedes every one of these without argument, which makes me trust the parts I have not checked more than I usually would.
\end{knote}\subsection*{Triples}

\walkfig{t1-hero}{The Triples lede and the live Borromean canvas.}

\Stu{This is the tab the rest of the app exists to reach. Two languages, one invariant: on the left $\bar\mu(123)$ through surfaces, on the right the R\'edei symbol $[13, 61, 937]$ through Frobenius. The canvas is live in a modest sense --- the three circles are drawn from stored centres and radii, the six crossings are computed as circle-circle intersections, and the over-under data is a cyclic table: $K_1$ over $K_2$, $K_2$ over $K_3$, $K_3$ over $K_1$, at both crossings of each pair.}

\Kle{Your $K_1$ label is guillotined. The top half of it is off the canvas.}

\Stu{It is. The label baseline is the circle's topmost point minus eight pixels, and $K_1$'s top sits thirteen pixels below the canvas edge, so the glyph runs off. Nobody caught it because the eye supplies the letter.}

\Kle{Second thing, and this one is not cosmetic. Those are round circles. Three round circles in space cannot form the Borromean rings. That is a theorem --- Freedman and Skora.}

\Stu{Correct, and the app cites exactly that theorem --- one layer down. What is on screen is a diagram: a plane shadow plus over-under data. The shadow is allowed to be round; the link it presents is not round in $S^3$. If you press ``Toggle 3D view (ellipses)'' you get genuine ellipses, semi-axes $2$ and $1.1$, one in each coordinate plane, and a note under them citing Freedman--Skora, J. Differential Geom. 25 (1987), for precisely the reason you just gave.}

\Kle{Then the citation is in the wrong place. The claim that needs defending is on this screen. The defence is behind a button and a hover.}

\Stu{Agreed. That is a fair charge and I have no answer to it.}

\Kle{What does the ``Isolate'' dropdown do?}

\Stu{Fades one component. Choose ``fade $K_1$'' and the app writes across the top: the remaining pair is visibly unlinked, one ring entirely on top. It is the cheapest way to see that pairwise vanishing is not a computation, it is a picture --- and the whole tab is about what survives it.}

\begin{knote}
Round circles on the canvas; Freedman--Skora is quoted, but only inside the 3D panel. The 2D picture is a shadow --- fine --- but he should say that where the shadow is, not two clicks away. Label $K_1$ clipped. Ellipses in the 3D view: $a = 2$, $b = 1.1$, one per coordinate plane. Ask later whether the 3D view needs the network.
\end{knote}

\walkfig{t2-steps}{The two four-step sequences: topology at left, R\'edei at right.}

\Stu{Same four beats down both columns. Step 1, the pairwise level vanishes: on the left $\mathrm{lk}(K_1,K_2) = \mathrm{lk}(K_1,K_3) = \mathrm{lk}(K_2,K_3) = 0$, computed from the diagram's signed crossings; on the right $(13/61) = (13/937) = (61/937) = +1$, computed by Euler's criterion. Step 2, the certificate the vanishing leaves behind: the longitude on the left, the R\'edei element $\alpha = 23 + 6\sqrt{13}$ on the right, with the identity $23^2 - 13 \cdot 6^2 - 61 \cdot 1^2 = 529 - 468 - 61 = 0$ shown as an evaluated residual. Step 3, evaluate the certificate: Magnus expansion on the left, Tonelli--Shanks for $\sqrt{13} \bmod 937$ on the right --- it returns $\{386, 551\}$. Step 4, read the invariant.}

\Kle{Stop at step 1. How does the sign cancellation actually come out of the code?}

\Stu{For each pair the app finds the two intersection points of the two circles, takes the counterclockwise tangent at each point on each circle, decides which strand is over from the cyclic table, and sums the sign of the cross product of over with under. Divided by two, that is the linking number.}

\Kle{And for two overlapping circles with the same strand on top at both crossings, those two signs are automatically opposite. I can see that without the machine --- the tangent frame reverses between the two intersection points. So step 1 is honest but it is not deep.}

\Stu{It is not deep. It is arithmetic on a picture.}

\Kle{Now step 2. You write ``that it is the specific commutator $\ell_3 = [x_1, x_2]$ is Milnor's Wirtinger computation, taken here as the model word rather than recomputed from the drawing.'' So the diagram gives you the vanishing, and the non-vanishing --- the thing the tab is named for --- is typed in.}

\Stu{Yes. That is the central gap of this tab and I will not dress it up. Vanishing $\mathrm{lk}$ tells you the longitude lies in the commutator subgroup $[F, F]$ and nothing more --- the tooltip is explicit that zero exponent sums do not make a word a single commutator; it gives $[x_1,x_2]^2$ as the counterexample, by Wicks. Which commutator it is comes from Milnor's Wirtinger computation, cited, not recomputed. What is computed live is everything downstream: substitute $x_i \mapsto 1 + X_i$, invert the series, multiply out, and read the coefficient of $X_1 X_2$.}

\Kle{So ``computed live'' means computed live from a hand-supplied input.}

\Stu{For the model word, yes. The same standard the Quadruples tab applies to $B_4$. The app says so in the step itself rather than in a footnote, which is the most I can claim for it.}

\Kle{Third thing, and it is a presentation failure. Your left column's step 2 is five times the height of the right column's step 2, so by the bottom of the screen your step 4 is opposite blank space and your step 3 is opposite their step 4. You have built a mirror and then let the halves slide.}

\Stu{The two columns are independent flows in a two-column grid; nothing ties row $k$ to row $k$. It is exactly the defect you name.}

\begin{knote}
Structure of the tab is right: vanish, certificate, evaluate, read --- twice. What is actually computed: crossings, Legendre symbols, the R\'edei identity residual, Tonelli--Shanks ($\sqrt{13} \bmod 937 = \{386, 551\}$), the Magnus coefficient. What is not: the longitude word itself. That is a fixture and he says so on the step. Mirror rows drift out of alignment --- fix with a row grid.
\end{knote}

\walkfig{t3-verdicts}{The two verdict banners: $\bar\mu(123) = 1$ and $[13, 61, 937] = -1$.}

\Kle{One banner says $1$. The other says $-1$. You are asking me to read those as the same statement.}

\Stu{I am, and the app hedges it in exactly one place --- the ``(sign?)'' tooltip on the left banner. The honest content is: $\bar\mu(123) = \pm 1 \neq 0$, mirroring $[13,61,937] = -1 \neq +1$. Not $1 = -1$.}

\Kle{Then the value groups are different and you should lead with that, not hide it in a hover. $\bar\mu$ is an integer. The R\'edei symbol lives in $\{\pm 1\}$.}

\Stu{The Q\&A tab has a whole entry on it, and it draws the line where you just drew it: the knot side has the full integers --- $\mathrm{lk} \in \mathbb{Z}$, $\bar\mu \in \mathbb{Z}$, and the Magnus coefficients are computed over $\mathbb{Z}$ before any reduction --- while $\mathrm{Gal}(\mathbb{Q}(\sqrt{p})/\mathbb{Q})$ has order $2$, so the Legendre symbol is genuinely a mod-$2$ linking number. The matching is $\bar\mu(123) \bmod 2 = 1$ against the nontrivial value of the symbol.}

\Kle{And the sign on the left is convention.}

\Stu{Orientation and component ordering. Reversing one component negates it; the mirror word $[x_2, x_1]$ gives $-1$. The app says the $+1$ is read off the model word, not derived from the drawn diagram --- which would need the Wirtinger presentation it does not compute.}

\Kle{Fine. Then the banner should read ``$\bar\mu(123) = \pm 1$, nonzero'' and the tab would lose nothing and gain a reader who is not me.}

\walkfig{t4-towers}{The cover tower --- mod-$2$ Heisenberg, deck $D_4$, $8$ sheets --- opposite the field tower $k_{13,61}$.}

\Stu{Left, a purely topological tower: the base $X = S^3 \setminus (K_1 \cup K_2)$, two double covers $X_1$ and $X_2$ --- the covers the linking numbers probe --- their join $X_{12}$, the $(\mathbb{Z}/2)^2$ cover, and on top $\tilde{X}$, the mod-$2$ Heisenberg cover, deck group $N_3(\mathbb{F}_2) \cong D_4$, $8$ sheets. Right, the field tower: $\mathbb{Q}$, $\mathbb{Q}(\sqrt{13})$, $\mathbb{Q}(\sqrt{61})$, $\mathbb{Q}(\sqrt{13},\sqrt{61})$, and $k_{13,61} = \mathbb{Q}(\sqrt{13},\sqrt{61},\sqrt{\alpha})$. The edge from the join to the top is labelled on both sides: two sheets where $\gamma$ unwinds, degree $2$ via $\sqrt{\alpha}$.}

\Kle{The caption underneath says the lattice of intermediate covers mirrors the subgroup lattice of the deck group. It does not. $D_4$ has three subgroups of index $2$, so there are three quadratic subfields of $k_{13,61}$. You have drawn two. Where is $\mathbb{Q}(\sqrt{793})$?}

\Stu{Missing, and so is the corresponding diagonal double cover on the left. The drawing is a sublattice, not the lattice: five nodes against $D_4$'s ten subgroups. The two I show are the two the linking numbers probe, which is the tab's narrative, but the caption claims the general correspondence and the picture does not deliver it.}

\Kle{The correspondence itself is true --- I am not disputing Galois theory. I am disputing a caption that promises a lattice and delivers a chain with a fork in it. Say ``the part of the lattice this computation uses'' and you are clean.}

\Stu{Noted, and it is a one-word fix.}

\Kle{One more. You promise ``box for box'' twice, above both towers. Look at the bottoms. $\mathbb{Q}$ sits a good deal higher than $X = S^3 \setminus (K_1 \cup K_2)$.}

\Stu{The two figures have different viewBox heights, $216$ against $190$, and no shared baseline. The top boxes happen to land level and the drift accumulates downward. The claim in the prose is stronger than the geometry on screen.}

\walkfig{t5-tower-caps}{The tower captions: $4$ loops winding twice against $937$'s $4$ primes of residue degree $2$.}

\Stu{This is the paragraph I would keep if I had to throw away the rest of the tab. On the left, $\rho(K_3)$ in $N_3(\mathbb{F}_2)$ has entries $(\mathrm{lk}, \mathrm{lk}, \bar\mu) \bmod 2 = (0, 0, 1)$; the app powers that matrix literally until it hits the identity and gets order $2$; so $K_3$'s preimage upstairs is $8/2 = 4$ loops, each winding twice. On the right, Frobenius at $937$ has order $2$, so there are $4$ primes above $937$ of residue degree $2$. Then the line I care about: loops times winding equals $4 \times 2 = 8 = |\mathrm{deck}|$; and $g \cdot f = 4 \times 2 = 8 = [k_{13,61} : \mathbb{Q}]$. The same $|G| = \mathrm{count} \times \mathrm{size}$ bookkeeping, twice.}

\Kle{That is the good sentence, and I will say so plainly: that is the only place so far where the two sides are not analogy but the same statement about a finite group acting on a fibre. Nothing is being smuggled.}

\Stu{Thank you.}

\Kle{Now let me try to break it. Both orders --- the matrix order on the left, the Frobenius order on the right --- are computed by the same subroutine. You call one function, feed it two inputs, and then present the agreement as a mirror. That is not a mirror. That is one calculation printed twice.}

\Stu{That is the sharpest form of the objection and it is half right. The shared routine takes a triple of bits and returns the order of the corresponding unitriangular matrix. What differs is what produces the bits. On the left, the corner bit is $\bar\mu(123) \bmod 2$ out of the Magnus expansion. On the right, the corner bit is the Legendre symbol of $\alpha = 23 + 6\sqrt{13}$ reduced mod $937$, after Tonelli--Shanks. Two entirely disjoint computations arriving at the same bit; the group theory downstream is deliberately shared, because the assertion is precisely that the bookkeeping is the same on both sides.}

\Kle{Then the mirror lives in the bit, not in the order.}

\Stu{Yes. And the app never claims otherwise --- but it also never says it that crisply, and it should.}

\Kle{And the factor twins at the end, $(1 - t^2)^{-4}$ against $(1 - 937^{-2s})^{-4}$?}

\Stu{Both assembled from the order and the degree. That is the Euler factor of $\zeta_{k_{13,61}}$ at $937$ on the right, and the corresponding orbit factor on the left, and the Zeta tab works the product through.}

\begin{knote}
The $|G| = \mathrm{count} \times \mathrm{size}$ paragraph is the best thing in the app. $4$ loops $\times$ winding $2 = 8$ sheets; $g = 4$, $f = 2$, $gf = 8 = [k_{13,61} : \mathbb{Q}]$. Real, not analogical. Pressed him on the shared order routine --- he conceded correctly: the independence is in the corner bit (Magnus versus Legendre-of-$\alpha$), not in the group theory. Tower caption oversells the lattice; drawing shows $5$ nodes of $D_4$'s $10$ subgroups, $\mathbb{Q}(\sqrt{793})$ absent.
\end{knote}

\walkfig{t6-whymu}{The why-box defining $\bar\mu(123)$ by tubing and triple points.}

\Stu{This is the geometric definition. Because every pairwise $\mathrm{lk}$ vanishes, each Seifert surface $\Sigma_i$ can be chosen disjoint from the other two components --- punctures come in cancelling pairs, tube them away. Then $\Sigma_1 \cap \Sigma_2$ has no boundary at all: it survives as a closed curve. And its intersection with $\Sigma_3$, rather than its boundary, becomes the well-defined count: $\bar\mu(123)$ is the signed sum over points of $(\Sigma_1 \cap \Sigma_2) \cap \Sigma_3$, independent of all choices, by Milnor --- because the pairwise linking vanished first. For the Borromean rings the three surfaces meet in a single triple point.}

\Kle{This is the definition of the invariant the entire tab computes, and it is collapsed behind a button labelled in small grey type. Why is it not the first paragraph?}

\Stu{Because the tab was built to lead with the two columns. It is the wrong ordering and I do not have a defence.}

\Kle{Now the mathematics. ``Tube them away'' is doing a great deal of work. Tubing changes the surface --- genus goes up, the surface is no longer the one you drew. Why does the triple-point count not move when I tube differently?}

\Stu{Independence is Milnor's, cited here, not proved in the app. The sketch the app does carry is in the triple-point tooltip on rung $3$: triple points are exactly where the closed curve $\gamma = \Sigma_1 \cap \Sigma_2$ meets $\Sigma_3$, so their signed count is the intersection number $\gamma \cdot \Sigma_3 = \mathrm{lk}(\gamma, K_3)$; trading one surface for another changes the count by an intersection with a closed surface, and that vanishes because the pairwise linking already did.}

\Kle{That is the right sketch. It is also the argument I would have given, so I believe the app read Milnor rather than paraphrasing a survey. But notice what has happened: the geometric definition is a story, and the number on the banner came from Magnus. Nowhere does this app count a triple point.}

\Stu{Nowhere. No surface is ever constructed in code. The surfaces are pictures; the invariant comes from the free group.}

\walkfig{t7-cockpit}{The certificate cockpit, seeded with $(5, 29, 109)$ and certificate $(7, 2, 1)$.}

\Stu{Honesty first, and the app leads with it: the pair $(13, 61)$ upstairs is pinned because the app ships R\'edei's normalized certificate $(23, 6, 1)$ as data. It verifies certificates; it does not search for them. This cockpit takes yours: primes $p_1, p_2, q$ and a certificate $(x, y, z)$ with $x^2 - p_1 y^2 - p_2 z^2 = 0$. The engine runs every definedness and normalization check and names the one that fails.}

\Kle{Name them.}

\Stu{Primality of all three; the congruence $p_1 \equiv p_2 \equiv q \equiv 1 \pmod 4$; the identity itself, reported as a residual; primitivity, $\gcd(x,y,z) = 1$; $y$ even; $x - y \equiv 1 \pmod 4$; the three Legendre conditions; and a check that $q$ does not divide a conjugate of $\alpha$.}

\Kle{Why does primitivity matter? If $(x,y,z)$ solves the form, so does $(kx, ky, kz)$.}

\Stu{Because scaling the certificate by $k$ twists the symbol by $(k/q)$ --- the field $\mathbb{Q}(\sqrt{13},\sqrt{61},\sqrt{k\alpha})$ is genuinely a different field. R\'edei's normalization is what makes the symbol independent of which solution you found. There is a self-test that feeds $(21, 6, 3)$ for $(5,29,109)$ --- the same solution scaled by three --- and asserts the engine rejects it on primitivity.}

\Kle{Good. Then the rejection taxonomy is the teaching content and the verdict is decoration.}

\Stu{That is the app's own framing, in that sentence.}

\walkfig{t8-cockpit-run}{The cockpit verdict for the seeded triple.}

\Stu{Run it and you get $[5, 29, 109] = +1$: $q$ splits completely in $k_{5,29}$, conjugates $49$ and $74$ mod $109$, both Legendre $+1$. Every check passed --- identity, primitivity, R\'edei normalization, definedness, primality.}

\Kle{So your fresh, bring-your-own example is unlinked. The one triple a visitor is handed at the door is the boring one.}

\Stu{Yes. The seeded triple demonstrates that all five gates open; it does not demonstrate a triple linking. The only $-1$ in the app is the pinned $(13, 61, 937)$.}

\Kle{That is worse than it sounds. A reader who never scrolls up sees a machine that says $+1$, and the tab's whole thesis is about the case where the answer is $-1$. Seed it with a linked triple and put the split one one click away, or label this one ``definedness demo'' in the box, not in the paragraph above it.}

\Stu{Both would be improvements. The paragraph does say it verifies rather than searches, but it does not say ``and the example you are looking at is the trivial value.''}

\Kle{Does the engine ever produce a $-1$ for a triple I invent?}

\Stu{It will, if you supply a normalized certificate. It will not find one for you. That is the honest boundary of the tool: no search, only verification.}

\begin{knote}
Cockpit: verification only, no search --- correctly advertised. Checks: primality, all three $\equiv 1 \bmod 4$, the form identity, $\gcd = 1$, $y$ even, $x - y \equiv 1 \bmod 4$, three Legendre conditions, $q \nmid$ conjugate. Primitivity matters because scaling by $k$ twists by $(k/q)$ --- self-tested with $(21,6,3)$. But: seeded example returns $+1$. The demo case is the case the tab is not about. Reseed.
\end{knote}

\walkfig{t9-ladder-intro}{The section head ``Climbing the ladder with surfaces --- what the zeta counts''.}

\walkbeat{He scrolls down past the head to the paragraph under it.}

\Stu{The columns above say what the invariants are; the ladder says what they count. Left rail, the surfaces of the Borromean rings rung by rung; right rail, the same climb for $(13, 61, 937)$. Each rung closes with the count its zeta function records --- sheets and loops upstairs on the left, primes above and residue degrees on the right. And the disclaimer, in the paragraph: the figures are schematic, honest cartoons of the general-position picture, not computed embeddings, while every number on the ladder is computed live.}

\Kle{``Honest cartoons'' is a good phrase and I will steal it. But you have just given me two categories --- schematic figure, live number --- and I am going to hold you to the boundary between them on every rung.}

\Stu{You should. There is exactly one figure in this section that is computed, and it is labelled so in bold.}

\walkfig{t10-rung1}{Rung $1$: $K_2$ piercing $\Sigma_1$ twice with opposite signs, against $(13/61) = +1$.}

\Stu{Rung $1$ is the pairwise count. A Seifert surface $\Sigma_1$ spans $K_1$; $K_2$ crosses it twice, once each way, dashed below the surface, plus and minus at the two punctures, and the caption says ``net $0$''. On the right, $(13/61) = +1$ computed by Euler's criterion, so $61$ splits in $\mathbb{Q}(\sqrt{13})$: two primes upstairs, each fixed by Frobenius. Then the counted line: in the double cover of $S^3 \setminus K_1$, $K_2$'s preimage is $2$ loops, from the parity of $\mathrm{lk}$; on the right, $2$ primes above $61$ of residue degree $1$, Euler factor $(1 - 61^{-s})^{-2}$ of $\zeta_{\mathbb{Q}(\sqrt{13})}$. Split equals two sheets equals two fixed points.}

\Kle{The sentence in the middle of your left panel does the work: ``computed from the Borromean diagram at the top of this tab, its signed crossings, not from this schematic.'' So the surface is a picture and the number is crossings.}

\Stu{Yes. No Seifert surface exists anywhere in this codebase. The surface language is how the invariant is explained; the arithmetic on the diagram is how it is produced.}

\Kle{I want that stated once, loudly, rather than five times quietly in figure captions. A manifold theorist reading this will assume you built a spanning surface, because that is what the pictures show.}

\Stu{That is the right criticism and it is the one I would act on first.}

\Kle{Second, and this is going to sound petty until you look at it. Your right rail reads ``$(13/61) =$'' on one line, then ``$+1$'' alone on the next line, then a line starting with a comma. That is not a line wrap. That is broken.}

\Stu{It is a CSS bug and I can tell you exactly what it is. The arithmetic rail is a flex container with column direction, so every child element becomes its own flex item on its own line. The bold $+1$ is a child element. So is the tooltip span further down. Every emphasized value in all three arithmetic rails gets torn onto its own line, with the surrounding prose orphaned around it.}

\Kle{So it repeats on rung $2$ and rung $3$.}

\Stu{On every rung, and on the tooltip in rung $2$ as well --- that is why ``the cover where $\gamma$ unwinds'' sits alone with a colon dangling under it.}

\Kle{Then it is one line of stylesheet, and it is disfiguring the half of the ladder you most want a number theorist to read.}

\begin{knote}
Rung $1$ honest: $\mathrm{lk} = 0$ from crossing signs, not from any surface. No Seifert surface is ever built in code --- say it once at the top, not in captions. Arithmetic rail is a flex column, so every $\langle b \rangle$ becomes its own line: ``$(13/61) =$'' / ``$+1$'' / ``, computed live''. One stylesheet line. Fix before anyone else sees this.
\end{knote}

\walkfig{t11-rung2}{Rung $2$: the three-step storyboard, the R\'edei certificate arithmetic, and the counted line.}

\Stu{The storyboard is three panels. One: the two spanning surfaces apart --- a flat disk with boundary $K_1$, a bowl with rim $K_2$. Two: the bowl pushed down through the disk; where the bowl's sides pass through, they cross it in exactly one closed circle, drawn bold, dashed where it is hidden below the disk. Three: that circle extracted and shown alone --- $\gamma = \Sigma_1 \cap \Sigma_2$, closed because rung $1$'s punctures cancelled, and not a component of the link.}

\Kle{Why one circle and not two? If I push a bowl through a plane I generically get a circle, yes, but I want to hear why this configuration gives one and not a pair.}

\Stu{The bowl meets the plane of the disk in a single closed curve because the bowl is a disk with one rim and the rim sits above the plane --- the intersection of the bowl's interior with the plane is the boundary of the region of the bowl lying below it, one component. If the bowl dipped through twice you would get two, and then $\gamma$ would be two circles and the count would run over both. The app does not compute this. It draws it.}

\Kle{Fair. Now the right rail. You match $\gamma$ to $\alpha = 23 + 6\sqrt{13}$ and then your headline is that the norm is $61$: $23^2 - 13 \cdot 6^2 = 529 - 468 = 61$, the second prime itself. That is a pretty coincidence. Why is it the right mirror?}

\Stu{It is not a coincidence and the R\'edei-field tooltip has the reason. The identity is $x^2 - p_1 y^2 = p_2 z^2$, so $\alpha \bar\alpha = p_2 z^2$, and here $z = 1$, so the norm is $61$ on the nose. That relation is the load-bearing one: it gives $\sqrt{\bar\alpha} = z\sqrt{p_2}/\sqrt{\alpha}$, which already lies in the field --- that is why the extension is Galois --- and the same relation is what selects $D_4$ over the quaternion group. So the norm being $p_2$ is exactly the statement that the top of the tower closes up, which is the arithmetic form of ``$\gamma$ is a closed curve''.}

\Kle{Good. That is a real answer and it is in the app rather than in your head?}

\Stu{In the tooltip on ``R\'edei field'', with R\'edei 1939 cited. The congruence conditions are attributed there too --- they are what kills ramification at $2$ and at infinity.}

\Kle{The counted line at the bottom says ``What is being counted: nothing yet.'' Explain that.}

\Stu{Rung $2$ builds the counting space rather than counting: the mod-$2$ Heisenberg cover of $S^3 \setminus (K_1 \cup K_2)$, $2^3 = 8$ sheets, deck group $N_3(\mathbb{F}_2) \cong D_4$, existing precisely because $\mathrm{lk} \equiv 0$ --- Morishita chapters $4$ and $8$, cited --- opposite the degree-$8$ field $k_{13,61}$. Rung $3$ counts orbits in it.}

\Kle{I like that a rung is allowed to count nothing. Most expositions would have invented a number here.}

\walkfig{t12-gamma-computed}{The one computed figure: standard ellipse position, flat disks, intersection segment sampled numerically.}

\Stu{This is the exception to ``schematic''. Standard ellipse position, the same coordinates as the 3D view: $K_1$ is the ellipse in $z = 0$ with semi-axes $2$ and $1.1$, $K_2$ the same ellipse in $x = 0$, and $\Sigma_1$, $\Sigma_2$ are the flat spanning disks. The app samples the line $x = z = 0$ at two hundred points, keeps those inside both disks, and reports $\Sigma_1 \cap \Sigma_2 = y \in [-1.10, 1.10]$. Separately it computes where $K_1$ punctures $\Sigma_2$ --- at $t = \pi/2$ and $3\pi/2$, with signs from the sign of $dx/dt$ --- and gets $y = \pm 1.1$ with opposite signs. The endpoints of the segment land exactly on the two cancelling punctures. That coincidence is asserted as a self-test to nine decimal places.}

\Kle{Then your one computed picture shows a segment. You have spent the last two rungs telling me $\gamma$ is a closed curve.}

\Stu{It shows the picture before the tubing, and the caption says so in those words: this figure computes the before, the cartoon above draws the after. The flat disks are not the surfaces the definition wants --- they still hit the other components. Tubing them away is what closes the arc into a circle, and the app does not compute the tubing.}

\Kle{So the honest reading is: the general-position intersection of the naive surfaces is an arc whose ends are precisely the two punctures that cancel, and the cancellation is what lets you close it up.}

\Stu{That is exactly the reading, and it is why this figure is worth having: it exhibits the mechanism --- ends of the arc equals the cancelling pair --- rather than the conclusion.}

\Kle{Second complaint, visual. Those two curves look nested. A blue ellipse sitting inside a green one. Nothing about that picture says two unlinked circles in general position, it says one loop inside another.}

\Stu{Projection artifact of the oblique view. And there is a real fact behind it: they genuinely are unlinked, because $K_2$ meets the plane $z = 0$ at $(0, \pm 2, 0)$, which is outside the disk $\Sigma_1$ --- $4/1.21 > 1$. So the picture is telling the truth about the linking and lying about the configuration.}

\Kle{And $K_3$ is not in this figure at all.}

\Stu{Not drawn. This figure is only about $\Sigma_1 \cap \Sigma_2$. $K_3$ appears on rung $3$, schematically.}

\Kle{I will say this: it is the only figure on the tab that could have been wrong and is not. Everything else is a drawing that cannot be checked. This one can be, and the endpoints agreeing with the punctures is a genuine test.}

\walkfig{t13-rung3}{Rung $3$: $K_3$ threading $\gamma$ once, $\bar\mu = \mathrm{lk}(\gamma, K_3)$, and $[13,61,937] = -1$.}

\Stu{$K_3$ passes through the circle $\gamma$ exactly once --- one triple point of $\Sigma_1$, $\Sigma_2$, $\Sigma_3$. The reading $\bar\mu(123) = \mathrm{lk}(\gamma, K_3) =$ signed count of triple points is cited: Cochran, Mem. AMS 84 (1990); Mellor--Melvin, Alg. Geom. Topol. 3 (2003); mod $2$, Morishita chapter $4$. The number itself is the Magnus coefficient from step $4$ above. On the right, $[13, 61, 937] = -1$: Frobenius at $937$ moves $\sqrt{\alpha}$, exactly as $K_3$ refuses to pull off $\gamma$. And the counted line: $\rho(K_3)$ has order $2$, so $K_3$'s preimage in the eight-sheet cover is $4$ loops each winding twice; on the right, $4$ primes above $937$ of residue degree $2$, Euler factor $(1 - 937^{-2s})^{-4}$ of $\zeta_{k_{13,61}}$.}

\Kle{Look at your own figure. ``surface spanning $\gamma$, drawn flat'' has the purple ellipse drawn straight through it. The words are struck out by the curve. And up at the top, ``the one pass through the surface:'' runs into the orange $K_3$.}

\Stu{Both true. The label sits at a point where the lower arc of the ellipse passes within a couple of pixels, and the leader text is long enough to reach $K_3$'s left edge. Neither is a rendering accident --- they are fixed coordinates in the SVG that were never checked against the curves.}

\Kle{Now the substantive one. Inside that figure you have written ``$\mathrm{lk}(\gamma, K_3) = 1$''. Is that computed?}

\Stu{No. It is a hardcoded text element in the SVG.}

\Kle{And the top of this tab says, in its navigation line, ``every number computed live''. And the paragraph introducing this ladder says ``every number on the ladder is computed live''.}

\Stu{Then the two statements are in tension and you have found the seam. The intended scope is: the numbers in the prose and the counted lines are computed; the figures are schematic. But a number printed inside a schematic figure is still a number on the ladder, and the blanket claim does not carve it out.}

\Kle{It is not a lie --- $\bar\mu(123) = 1$ is computed, twelve lines below, and the two agree. But the app's entire credibility rests on the distinction between computed and cited, and here it has quietly printed a third category: drawn.}

\Stu{Drawn. That is the right word and the app does not have it.}

\Kle{One more, mathematical. $\mathrm{lk}(\gamma, K_3)$ needs $\gamma$ oriented. Where does the orientation come from?}

\Stu{From the orientations of $\Sigma_1$ and $\Sigma_2$ --- the intersection orientation of the two surfaces. Reverse either surface and $\gamma$ reverses, and the sign of $\bar\mu$ flips, which is the same convention dependence the banner flags. The app never says this. It is a genuine omission: the sign discussion is entirely about component ordering and mirror words, and never about how $\gamma$ inherits a direction.}

\Kle{Add a sentence. It is the kind of thing that costs nothing and buys a reader's trust.}

\begin{knote}
Rung $3$: identification $\bar\mu = \mathrm{lk}(\gamma, K_3) = $ signed triple points, cited (Cochran 1990; Mellor--Melvin 2003; Morishita ch.\ 4 mod $2$); number from Magnus. Counted line correct: order $2$, $4$ loops winding twice, $(1 - 937^{-2s})^{-4}$. Caught: ``$\mathrm{lk}(\gamma, K_3) = 1$'' is drawn, not computed, inside a tab claiming every number is computed. Third category needed: computed / cited / drawn. Also: $\gamma$'s orientation is never explained --- it comes from the intersection orientation of $\Sigma_1$, $\Sigma_2$. Figure labels collide with the curves.
\end{knote}

\walkfig{t14-polys}{The $\Delta_{\mathrm{Bor}}$ coefficient grid: every lower coefficient vanishes, the top one is $\bar\mu$.}

\Stu{$\Delta_{\mathrm{Bor}}(t_1,t_2,t_3) = (t_1-1)(t_2-1)(t_3-1)$. Substitute $t_i = 1 + x_i$ and expand live, and you get the grid: constant $0$, all three linear coefficients $0$, all three quadratic coefficients $0$, and the coefficient of $x_1x_2x_3$ equal to $1$, highlighted, labelled $\bar\mu(123)$. Every coefficient below top degree vanishes --- mirroring the vanished pairwise linking --- and the single surviving coefficient is the triple invariant.}

\Kle{Is $\Delta_{\mathrm{Bor}}$ computed?}

\Stu{No. It is cited --- a standard Fox-calculus result, Milnor 1954 for the Borromean case, tables in Kawauchi. The app has no Fox calculus. What is computed is the expansion you are looking at, plus a Torres collapse to $(t-1)^4$ checked against the fibre's $b_1 = 4$ on the Zeta tab.}

\Kle{Then be careful about what the grid proves. You substitute $t_i = 1 + x_i$ into a product of three factors $(t_i - 1)$. That expansion is $x_1x_2x_3$ and nothing else. The grid is not a discovery, it is the distributive law confirming that seven zeros are seven zeros.}

\Stu{That is true. The grid is a display, not a test.}

\Kle{And the interesting claim is the one underneath: that the surviving top coefficient equals $\bar\mu(123)$, computed independently. Except it is not independent, is it? Your $\bar\mu$ came from the Magnus expansion of $[x_1, x_2]$, which you told me is Milnor's model word for this link. And $\Delta_{\mathrm{Bor}}$ is Milnor 1954. Both sides descend from the same paper.}

\Stu{They do. The cross-check has a common ancestor, and I will not claim it is two independent witnesses. What it does check is internal consistency: given Milnor's word and Milnor's polynomial, the app's own arithmetic on both agrees where the theory says it must. If I had made a sign error in the Magnus inversion, the grid would disagree.}

\Kle{That is a test of your code, not of the mathematics. Which is worth having --- I would rather have code I trust --- but the tab presents it as a mathematical cross-check, and the tooltip calls it ``two independent cross-checks''.}

\Stu{Then the word ``independent'' is doing work it has not earned. The Torres collapse against $b_1 = 4$ is genuinely a second constraint; the coefficient statement against $\bar\mu$ is not.}

\Kle{Correct. Downgrade that word and the sentence is true.}

\Stu{The right column, for completeness, says the matching statement --- and hedges harder than the left. The R\'edei symbol is the triple Massey product in Galois cohomology, defined because the cup products vanish, and equal to the arithmetic triple Milnor invariant by Morishita, Compositio Math.\ 140 (2004). But the analogous claim for the arithmetic Alexander series $f_S$ --- that its leading surviving coefficient reports the R\'edei symbol --- is offered as analogy, not theorem, because over $\mathbb{Q}$ the deck group is finite.}

\Kle{That is the correct hedge and I checked the reason on the Dictionary tab: no $\mathbb{Z}^r$ tower over $\mathbb{Q}$ unramified only at your primes. Good. The tab tells the truth about which of its two halves is a theorem.}

\begin{knote}
Grid is the distributive law, not a test --- he conceded. $\Delta_{\mathrm{Bor}}$ cited (Milnor 1954), no Fox calculus in the app. The ``two independent cross-checks'' claim: the coefficient-versus-$\bar\mu$ check shares its ancestor with the polynomial, so it tests his code, not the mathematics. The Torres collapse against $b_1 = 4$ is the real second constraint. Right column hedges correctly: Morishita 2004 is a theorem, the $f_S$ statement is analogy, and the reason ($f_S$'s deck group finite over $\mathbb{Q}$) is computed on the Dictionary tab.
\end{knote}

\subsection*{Ladder}

\walkfig{l1-intro}{The Ladder tab heading and its one-sentence intro.}

\Stu{The unipotent ladder. Each level of linking lives in a unipotent matrix group over $\mathbb{F}_2$; the superdiagonal carries the characters, which is the pairwise level, and each deeper diagonal is defined only when everything above it vanishes.}

\Kle{Whose linking?}

\Stu{Both, and the sentence does not say so --- that is the thinnest opening in the app.}

\Kle{And by what map? A matrix group does not contain linking numbers. Something has to map into it. On the Triples tab you had $\rho(K_3)$, so there is a representation somewhere, but it is not named here.}

\Stu{It is named in the cover caption on Triples: $\pi_1$ of the link complement surjects onto $N_3(\mathbb{F}_2)$, meridians going to the generator matrices, and the eight-sheet cover is the one that classifies. On the arithmetic side it is $\rho$ on Frobenius. Neither appears on this tab's opening, and it should --- this is the tab where the object is the group.}

\Kle{And ``characters''. To me a character is a homomorphism to $\mathbb{Z}/2$. Is that what you mean?}

\Stu{Yes --- the mod-$2$ linking number on one side and the Legendre symbol on the other, both as homomorphisms to $\mathbb{Z}/2$. The word is used unexplained here. The tooltip on ``unipotent'' explains the group --- upper triangular with ones on the diagonal, order $2^{n(n-1)/2}$, so $2$, $8$, $64$ for $n = 2, 3, 4$ --- but nothing explains ``character''.}

\walkfig{l2-matrices}{The $N_2$, $N_3$, $N_4$ unipotent matrix displays.}

\Stu{Three cards. $N_2(\mathbb{F}_2) \cong \mathbb{Z}/2$, order $2$, one entry $\chi$, highlighted, always defined. $N_3(\mathbb{F}_2) \cong D_4$, order $8$: superdiagonal $\chi_1$, $\chi_2$ and corner $r$. $N_4(\mathbb{F}_2)$, order $64$: superdiagonal $\chi_1$, $\chi_2$, $\chi_3$, second diagonal $r_{123}$, $r_{234}$, corner $\bar\mu$. The arrows between them are buttons: press the first and it lights the corner of $N_3$ and writes that the corner is defined because the previous level's linking numbers vanished; press the second and it lights all three deeper entries of $N_4$.}

\Kle{Your $N_3$ corner is called $r$ and your $N_4$ corner is called $\bar\mu$. One is R\'edei's letter and one is Milnor's. On the tab whose entire point is that the two sides are the same object, you have written the two sides in different alphabets, one card apart.}

\Stu{That is a real inconsistency and I have no story for it. The $N_4$ second diagonal is $r_{123}$, $r_{234}$ --- arithmetic --- and then the corner above them switches to topology. It should be one language with the other in the tooltip, and the tooltips do carry both.}

\Kle{Second. The three cards are not aligned. $N_2$'s superdiagonal sits lower than $N_3$'s, which sits lower than $N_4$'s. You are asking me to watch a corner move rightward and downward through three pictures, and the pictures themselves are sliding vertically.}

\Stu{The row is a flex container with items centred, so cards of different heights centre on a common midline instead of a common top. Top alignment would put all three matrices' first rows level and the growth would read cleanly.}

\Kle{Now a question. You say the second diagonal of $N_4$ is defined when the superdiagonal vanishes. Justify that the way you justified it for $N_3$.}

\Stu{Same argument one level up. Frobenius, or the meridian image, is only well defined up to conjugation. Conjugating a unitriangular matrix changes an entry on the $k$-th diagonal by products involving entries on shallower diagonals. So an entry is conjugation-invariant exactly when everything above it in that sense vanishes. That is what the tooltip on the ``why?'' link says for $N_3$, and it is the algebraic face of the whole app's slogan: the next invariant is defined because the previous level vanished.}

\Kle{Then the $N_4$ card is not a new theorem, it is the same lemma with more indices. Good --- that is what a ladder should be.}

\begin{knote}
$N_2, N_3, N_4$: orders $2$, $8$, $64$, $= 2^{n(n-1)/2}$. Definedness argument is uniform: conjugation shifts a deep entry by products of shallower ones, so an entry is invariant exactly when its predecessors vanish. Caught: $N_3$'s corner labelled $r$ (arithmetic), $N_4$'s corner labelled $\bar\mu$ (topological) --- two alphabets on one row, on the tab that exists to say they are one thing. Cards centre-aligned, so the diagonals do not line up across the arrows.
\end{knote}

\walkfig{l3-frob937}{The Frobenius matrix panel for $q = 937$.}

\Stu{For the pair $(13, 61)$, every odd prime $q$ not dividing $13 \cdot 61$ gets an actual matrix in $N_3(\mathbb{F}_2)$. The superdiagonal entries are its two Legendre symbols, zero for $+1$ and one for $-1$; the corner is its splitting in $k_{13,61}$. At $q = 937$: $(13/937) = (61/937) = +1$, so both superdiagonal entries are zero; then Tonelli--Shanks gives $\sqrt{13} \equiv 386$, so $\alpha \equiv 23 + 6 \cdot 386 \equiv 465 \pmod{937}$, and $465$ is a non-square mod $937$. The corner lights up. Matrix: ones on the diagonal, zeros on the superdiagonal, one in the corner.}

\walkbeat{Klein does the arithmetic on the back of a printout before answering.}

\Kle{$386^2 = 148996$, and $937 \times 159 = 148983$, difference $13$. Your square root is right.}

\Stu{Everything on that panel is computed at the press: two Euler-criterion symbols, and when both vanish, Tonelli--Shanks and one more Legendre symbol.}

\Kle{The matrix you have drawn is the central element of $D_4$ --- zeros on the superdiagonal, one in the corner. Its fixed field is the biquadratic field $\mathbb{Q}(\sqrt{13},\sqrt{61})$. So $937$ should split into four primes down there and each one should stay inert going up to $k_{13,61}$.}

\Stu{And that is what the thread underneath reports, from the other direction: order $2$ by literal powering of the matrix, so the eight-element cover falls into $4$ orbits of length $2$, local factor $(1 - 937^{-2s})^{-4}$. Four primes, residue degree two.}

\Kle{Consistent. And I notice the order is obtained by multiplying the matrix by itself until it is the identity, rather than by a rule.}

\Stu{Deliberately. There is a rule --- identity if everything vanishes, order $4$ if both $\chi$ are $1$, otherwise order $2$ --- and in this app the rule is the test and the powering is the computation, not the other way round.}

\Kle{That is the correct way round. I have seen too many programs that assert the classification and then test it against itself.}

\walkfig{l4-frob107}{The same panel for $q = 107$, and the two-tier domain note.}

\Stu{$q = 107$: $(13/107) = (61/107) = +1$, $\sqrt{13} \equiv 21$, so $\alpha \equiv 23 + 126 \equiv 42 \pmod{107}$, and $42$ is a square. So $107$ splits completely in the degree-$8$ field and Frobenius is the identity matrix: order $1$, eight orbits of length one, local factor $(1 - 107^{-s})^{-8}$.}

\Kle{$21^2 = 441 = 428 + 13$, and $428 = 4 \times 107$. Fine.}

\Stu{Now the caveat, and it is the most careful paragraph on the tab. $107 \equiv 3 \pmod 4$. The corner --- the splitting of $q$ in $k_{13,61}$, a field fixed once and for all --- is defined for every odd unramified $q$ once the superdiagonal vanishes, whatever $q$ is mod $4$. The R\'edei symbol $[13, 61, q]$ is the same number when moreover $q \equiv 1 \pmod 4$. That extra congruence buys the symbol its symmetry and reciprocity, which is what makes it a triple linking number matching $\bar\mu$ --- not the definedness of the corner. So the app's R\'edei routine rejects $q \equiv 3 \pmod 4$ while this matrix does not, and the self-tests pin both behaviours at $q = 107$ on purpose.}

\Kle{That is a genuinely careful distinction and most expositions would blur it. Here is my problem with it. On the topology side there is no congruence. $\bar\mu$ is defined for any link with vanishing pairwise linking, full stop. Your arithmetic side needs an extra hypothesis exactly at the point where you want to call the thing a linking number. So the mirror has a hole in it, and the hole is at the load-bearing spot. What is the topological shadow of $q \equiv 1 \pmod 4$?}

\Stu{The app's answer, on the Dictionary tab, is that $p \equiv 1 \pmod 4$ plays the role of orientability --- it tames ramification at $2$ and at infinity in the quadratic resolvents. That is the analogy it offers, and it is offered as an analogy. My honest answer to your question is that in this app the congruence has no independent topological content. It is a convention that makes the arithmetic side behave, and the app is careful never to claim a topological counterpart for it beyond the word ``orientability''.}

\Kle{Then say that in the domain note. You have written two hundred and fifty words of correct hedging and the reader has to assemble the punchline themselves.}

\Stu{And it is a single unbroken paragraph sitting above the buttons, so a reader arrives at the presets having skimmed it. The two-tier distinction is the second-most important sentence on the tab and it is in the middle of the block.}

\Kle{There is a worse consequence. A visitor presses $q = 107$, sees a matrix, sees ``splits completely'', and reads that as a R\'edei symbol value. Your engine refuses to compute a symbol there. The screen does not refuse anything.}

\Stu{The explanation text does carry the parenthetical --- for $107$ it says the splitting datum is well defined but the symmetric symbol asks for $q \equiv 1 \pmod 4$. But you are right that the matrix is the loud object and the caveat is the quiet one.}

\begin{knote}
$q = 937$: $\sqrt{13} \equiv 386$ (checked: $386^2 - 13 = 159 \cdot 937$), $\alpha \equiv 465$, non-square, corner $= 1$, central element of $D_4$, order $2$, $(1 - 937^{-2s})^{-4}$ --- consistent with $937$ splitting into $4$ primes in the biquadratic field and staying inert above. $q = 107$: $\sqrt{13} \equiv 21$, $\alpha \equiv 42$, square, identity, $(1 - 107^{-s})^{-8}$. Orders by literal powering, classification rule kept as the test --- correct discipline.

Two-tier domain: corner defined for all odd unramified $q$; the symmetric R\'edei symbol wants $q \equiv 1 \bmod 4$, which buys reciprocity, not definedness. Careful and correct. But: the topology side has no such congruence, so the mirror is asymmetric exactly at the linking-number claim, and he conceded there is no independent topological content --- ``orientability'' is the analogy, nothing more. Caveat buried in a $250$-word block above the buttons.
\end{knote}

\walkfig{l5-chebotarev}{The Chebotarev tally over all $427$ unramified odd primes below $3000$.}

\Stu{Press the button and the app runs the whole panel above for every odd prime under $3000$ except $13$ and $61$ --- Tonelli--Shanks included whenever both $\chi$ vanish --- and bins each one by its conjugacy class in $D_4$. Five classes: identity, corner, and three $\chi$-patterns. Counts $50$, $47$, $109$, $111$, $110$ out of $427$; observed $0.117$, $0.110$, $0.255$, $0.260$, $0.258$; predicted $1/8$, $1/8$, $1/4$, $1/4$, $1/4$. One millisecond.}

\Kle{Where do the predicted densities come from?}

\Stu{Chebotarev, cited, not proved: each class gets $|C|/|G|$. The identity and the central corner element are singleton classes, so $1/8$ each; the three $\chi$-patterns are classes of size two, so $1/4$ each. One plus one plus two plus two plus two is eight.}

\Kle{Now. Both of your singleton classes came in low --- $50$ and $47$ against an expected $53.4$ --- and all three doubletons came in high. Is that agreement or is that a deficit?}

\Stu{At this cutoff it is noise, and I can put a number on it, which the app cannot. The standard deviation for a singleton class is the square root of $427$ times $1/8$ times $7/8$, about $6.8$. So $50$ is about half a sigma low and $47$ about nine-tenths of a sigma low. Both ordinary.}

\Kle{And the apparent pattern --- two low, three high --- is forced, isn't it. The counts sum to $427$.}

\Stu{Exactly forced. The two singleton deficits are $3.4$ and $6.4$, total $9.75$; the three doubleton surpluses are $2.25$, $4.25$ and $3.25$, total $9.75$. Identically. There is no pattern to explain; the classes are negatively correlated because the total is fixed.}

\Kle{Good. Then here is the criticism, and it is the one I would put at the top of the list for this tab. Your table has no error bar. The caption says ``observed hugs predicted''. ``Hugs'' is a verb, not a measurement. A topologist reading this table has no way to tell whether $0.110$ against $0.125$ is a triumph or a problem --- and you just showed me it takes one line to answer.}

\Stu{There is no sigma anywhere in the app, on this table or any other. Adding one column would turn a verb into a number.}

\Kle{Do it. People believe numbers with error bars and distrust prose that asks them to feel reassured.}

\Stu{The caption is otherwise careful --- it says the equidistribution itself is Chebotarev's theorem, cited not proved, and that the app is a demonstration bench, not a proof. The Q\&A entry says the same thing more bluntly: verifying a density needs primes to infinity, which no browser computes.}

\Kle{That framing I have no quarrel with. Exhibiting instances and calling them instances is honest work. It is the one adjective --- ``hugs'' --- that is doing unearned persuasion.}

\walkbeat{He looks at the caption strip at the bottom of the tab.}

\Kle{Last thing. Your closing line reads $\mathrm{Gal}(k_{13,61}/\mathbb{Q}) \cong D_4 \subset N_3(\mathbb{F}_2)$. But your own card, two screens up, says $N_3(\mathbb{F}_2) \cong D_4$. If they are isomorphic, the inclusion sign is at best vacuous and at worst tells me the Galois group is a proper subgroup.}

\Stu{It is a leftover from the general statement --- for higher $n$ the image does sit properly inside $N_n(\mathbb{F}_2)$ --- but at $n = 3$ it is all of it, and the strip should say $\cong$. Not false, but it undercuts the sentence it is trying to land.}

\Kle{Then it is the last thing anybody reads on this tab, and it is the one that has to be exactly right.}

\begin{knote}
Tally: $427$ primes below $3000$, counts $50 / 47 / 109 / 111 / 110$ against $1/8, 1/8, 1/4, 1/4, 1/4$. Singleton sigma $\approx 6.8$; deficits are $0.5$ and $0.9$ sigma; the ``two low, three high'' pattern is forced by the fixed total ($9.75$ either way, exactly). He produced all of that on demand --- but the app produces none of it. No error bar anywhere. ``Observed hugs predicted'' is the only weak sentence in an otherwise scrupulous caption. Add a sigma column.

Closing strip writes $D_4 \subset N_3(\mathbb{F}_2)$ where the tab elsewhere asserts $\cong$. Last line on the tab; make it exact.

Overall on these two tabs: the machinery is real and the citations are placed where the theorems are. The recurring failure is not mathematical, it is that the honest qualification is always one layer below the confident claim --- in a tooltip, in a collapsed box, in the middle of a long paragraph. He knows this. He conceded every instance without argument, which is the reason I believe the parts I did not check.
\end{knote}\subsection*{Quadruples}

\walkfig{q1-intro}{The Quadruples opener: one line of scope-setting above the two columns.}

\Stu{One sentence, and it is the whole thesis of the tab. R\'edei's classical construction stops at triples; Amano's 2014 paper and the unipotent ladder from the tab before this one run one more level --- the same argument, one matrix size up. Everything underneath is that sentence cashed out twice: topology on the left, arithmetic on the right.}

\Kle{``The same argument, one matrix size up'' is either a very strong claim or an empty one. Which is it?}

\Stu{Strong in that the only thing that moves is the group: $N_3(\mathbb{F}_2)$ becomes $N_4(\mathbb{F}_2)$, and the shape --- an order rule in a unipotent group, a splitting condition, a cover --- is unchanged. Not empty, because Amano had to prove an independence theorem to make the four-term bracket a function of the four primes at all. The tab shows you exactly where that theorem gets used. It is step 3.}

\Kle{I will hold you to step 3.}

\walkfig{q2-b4}{The four-component Brunnian link, redrawn from the original path data, with the live pairwise-linking note beneath it.}

\Stu{This is Wikipedia's Brunnian-4loops figure by AnonMoos, public domain, redrawn from its exact path data --- the crossing gaps sit at computed intersections and the over/under comes from the original's layer order. The note underneath is live. It walks the diagram's signed crossings and reports $lk(K_1,K_2) = 0$, $lk(K_2,K_3) = 0$, $lk(K_2,K_4) = 0$, $lk(K_3,K_4) = 0$.}

\Kle{That is four numbers. A four-component link has six pairs. Where are the other two?}

\Stu{In the same sentence, stated differently: $K_1 \cap K_3 = K_1 \cap K_4 = \emptyset$. Green never meets red or blue anywhere in this drawing, so there are no crossings to sign. Those two pairs are unlinked because they are disjoint on the page, not because a signed sum came out zero. I would rather the note said that in two clauses than pretend six sums were computed.}

\Kle{Accepted --- and it is checkable by eye, which is the right kind of honesty. Now a complaint. Your $K_4$ label is written in blue directly across the blue strand. I read it as ``K4 (b1ue)'' for a second. Nudge it out into the white space; there is plenty at that corner.}

\Stu{Fair. The four labels are placed at fixed canvas coordinates and the blue one collides with its own curve. That is a five-pixel fix I have not made.}

\Kle{And the Brunnian property itself --- ``cut any loop and the other three fall apart'' --- you are not computing that.}

\Stu{No. The honesty note below the figure says it outright: the Brunnian claim is the figure's documented property, and the algebraic certificate $\bar\mu(1234) = 1$ is computed for Milnor's model $B_4$, not for the pictured diagram. Getting $\bar\mu$ off this drawing would need its Wirtinger presentation, which the app does not build.}

\Kle{So there are two objects on this screen and only one of them is computed. Say that out loud every time, because a reader who skims will merge them.}

\walkfig{q3-delete}{The delete-a-ring note: the algebraic shadow of cutting a component.}

\Stu{The interactive part. Pick a component, press Delete ring, and the drawing fades it and drops the pairs it was in. But the sentence here is about the algebra. Delete $K_1$, $K_2$, or $K_3$ --- algebraically, set that meridian $x_i = 1$ --- and the word $\ell_4 = [[x_1,x_2],x_3]$ collapses to $1$. That is $K_4$'s longitude in the model, written in the meridians of the other three.}

\Kle{Check it in front of me. $x_2 = 1$.}

\Stu{$[[x_1,1],x_3] = [1,x_3] = 1$. Same for $x_1$, and $x_3 = 1$ kills the outer bracket directly. All three collapse.}

\Kle{And $K_4$?}

\Stu{Different, and the note says so: $\ell_4$ is $K_4$'s own longitude, so there is no $x_4$ to set --- the word leaves with its component. Deleting $K_4$ is a real deletion in the picture but a vacuous one in this word.}

\Kle{Good. But notice what the note is careful about --- it says the collapse means ``the Milnor data vanish,'' not that the sublink is trivial. Those are not the same statement.}

\Stu{Right, and the note draws that line explicitly: the sublinks' actual triviality is the model's construction, not something read off the collapsing word. The word vanishing is the shadow, not the thing.}

\begin{knote}
Quads tab. Two $B_4$'s live here: the drawn diagram (linking numbers computed, Brunnian cited) and Milnor's model (longitude $[[x_1,x_2],x_3]$, $\bar\mu$ computed). App is honest about it but the two sit six inches apart on one screen. Label collision on $K_4$. Deletion demo: word collapse verified by hand, three cases, correct.
\end{knote}

\walkfig{q4-magnus}{The Magnus expansion box for the model longitude.}

\Stu{This is the computation. The Magnus expansion embeds the free group into power series in noncommuting variables, $x_i \mapsto 1 + X_i$. Expanding $[[x_1,x_2],x_3]$ gives $1 + X_1X_2X_3 - X_2X_1X_3 - X_3X_1X_2 + X_3X_2X_1$ plus degree four and up, and $\bar\mu(1234)$ is the coefficient of $X_1X_2X_3$, highlighted: $1$. About forty lines of truncated noncommutative series arithmetic, running at page load.}

\Kle{Milnor invariants of length four are only defined modulo an indeterminacy coming from the shorter ones. What is your indeterminacy here?}

\Stu{Zero, and for the reason the tab spent its first half establishing: every pairwise linking number vanishes and every three-component sublink is trivial in the model, so the ideal generated by the lower invariants is zero and the length-four invariant is an honest integer. That is why this rung exists at all --- the deeper diagonal is only defined once everything above it vanishes. It is the ladder's discipline, one step up.}

\Kle{Then the degree-three part had better be exactly those four words, with coefficients of absolute value one, or your indeterminacy story is decoration.}

\Stu{The self-test suite pins precisely that: the degree-at-most-three support is exactly the set $\{1, +X_1X_2X_3, -X_2X_1X_3, -X_3X_1X_2, +X_3X_2X_1\}$, all coefficients of modulus one, and it asserts the displayed text matches. If someone edits the expansion code, that test goes red.}

\Kle{That I like. Most demos of this kind print a number and dare you to disbelieve it.}

\walkfig{q5-steps}{The arithmetic column: five computed steps for Amano's quadruple $(5, 8081, 101, 449)$.}

\Stu{Amano's quadruple, computed live. Step 1: all six pairwise Legendre symbols, all $+1$ --- no two of these primes are linked. Step 2: three R\'edei elements with their identity and normalization checks, all residuals zero. Step 2b: five triple R\'edei symbols, all $+1$ --- the triple level is trivial too. Step 3 is the new level: $\theta = 25 + 2\sqrt{5} + 2\sqrt{101}$. Step 4 takes square roots mod $449$ by Tonelli--Shanks. Step 5 reads the four conjugates of $\theta$ mod $449$ and their Legendre symbols.}

\Kle{Stop at 2b for a moment. You list $[5, 8081, 101]$ and $[5, 101, 8081]$ separately. Those are the same three primes.}

\Stu{Deliberately. They are computed through two entirely different R\'edei elements --- $\alpha = 241 + 100\sqrt{5}$ for one, $11 + 2\sqrt{5}$ for the other --- and they agree. That is one instance of R\'edei reciprocity, witnessed rather than asserted. The tooltip is careful that this exhibits one transposition, not all of $S_3$, and that R\'edei's own symmetry argument from 1939 is regarded as incomplete; the complete proofs are Corsman 2007 and Stevenhagen 2018.}

\Kle{Now step 3, which you promised me. Why is that identity not enough?}

\Stu{Because $X^2 - 101Z^2 = \alpha$ does not pin $\theta$ down. Other solutions satisfy the identity and certify the opposite verdict. Concretely, scaling by a norm-one factor --- Hilbert 90, $\xi = 10 + \sqrt{101}$ has norm $-1$ --- gives a totally negative $\theta''$, hence an imaginary field ramified at infinity, and $449$ cannot see the difference because $(-1/449) = +1$. So the app checks total positivity: the four real embeddings come out near $49.57$, $9.37$, $40.63$, $0.43$, all positive.}

\Kle{Zero point four three. That is not comfortable.}

\Stu{It is not, and it is real: $25 - 2\sqrt{5} - 2\sqrt{101}$ is genuinely close to zero. The check is exact in spirit --- the identity residuals are BigInt and land on zero --- but the positivity test is floating point on four doubles. At this size that is a hundred-and-fifty-bit margin, so it is safe, but I would not call it certified. The exploit that the self-tests reject is caught by sign, not by magnitude, so it survives that criticism; a $\theta$ engineered to sit at $10^{-16}$ would not be caught.}

\Kle{Write that down as a known limitation rather than hoping nobody multiplies it out. And the second condition?}

\Stu{The 2-adic normalization, which is cited and not computed --- Amano's, and the tab says so in the step-3 tooltip and again in the banner. Positivity handles infinity; the prime $2$ is on trust.}

\Kle{One more, small. Step 5 prints your symbols as ``$\{-1, -1, -1, -1\}$'' with a typewriter hyphen, and the banner two inches below prints $-1$ with a proper minus. On a page this typographically fussy, that reads as a bug.}

\Stu{It is the same code path in both places emitting a raw JavaScript number in one and a formatted string in the other. You are right that it looks like the two disagree.}

\begin{knote}
Step 3 is the load-bearing one. Identity alone underdetermines $\theta$; Hilbert-90 twist gives a totally negative competitor invisible mod $449$ since $(-1/449) = +1$. Positivity kills it --- but smallest embedding is $0.43$ and the test is floating point. 2-adic condition cited, not computed. Reciprocity: one transposition exhibited, not $S_3$; app says so.
\end{knote}

\walkfig{q6-banner}{The quadruple-symbol verdict.}

\Stu{$[5, 8081, 101, 449] = -1$. No two linked, no three linked, all four linked: the arithmetic $B_4$. The banner is generated from step 5's four symbols --- if they had all come out $+1$ it would say $449$ splits completely and unlinked at the fourth level, and if they had disagreed it would refuse to assert anything.}

\Kle{Can they disagree?}

\Stu{No. At a prime modulus the four conjugate values are Frobenius-stable, so agreement is forced. That branch is an internal guard, and the code comments call it that rather than dressing it up as a mathematical contingency. It exists so that a future edit that breaks Tonelli--Shanks produces a refusal instead of a confident wrong answer.}

\Kle{Then the interesting content of this banner is the tooltip, not the banner. It is the tooltip that tells me the four-place notation is a theorem --- that the value depends on four primes and not on $\theta$.}

\Stu{Yes, and that is a fair criticism of the layout: the headline is one symbol and the justification is a hover away.}

\walkfig{q7-tower}{The 64-sheet cover tower with deck group $N_4(\mathbb{F}_2)$, and its computed caption.}

\Stu{The topological mirror, box for box against the field tower. The base is $X = S^3 \setminus (K_1 \cup K_2 \cup K_3)$; three meridian double covers; two four-fold covers through the shared $K_2$; two 8-sheet mod-2 Heisenberg covers $H_{12}$ and $H_{23}$ with deck group $D_4$; their join at 32 sheets; and on top the 64-sheet cover with deck $N_4(\mathbb{F}_2)$.}

\Kle{Why is $K_4$ not in the base?}

\Stu{Because it is the object being read, not part of what is being covered --- exactly as $449$ is the prime whose splitting is read in a field ramified only at $5$, $8081$, $101$. The map $\rho$ sends the meridians $x_1, x_2, x_3$ to the superdiagonal generators, and $K_4$'s class in $\pi_1(X)$ is its longitude $[[x_1,x_2],x_3]$, which lands on the corner. So $\rho(K_4)$ has entries all zero except $\bar\mu(1234) \bmod 2 = 1$ at the corner.}

\Kle{And the corner element has order two, so its preimage is thirty-two loops of double length. Fine. Say the counting rule.}

\Stu{Loops times winding equals the order of the deck group: $32 \times 2 = 64$. The caption states the arithmetic twin in the same breath --- $449$'s 32 orbits of size 2, conjugate symbols recomputed live --- and prints the two factors side by side: $(1 - t^2)^{-32}$ and $(1 - 449^{-2s})^{-32}$.}

\Kle{Now the label on the left edge of the tower says ``2 sheets, $\gamma_{12}$ unwinds.'' What is $\gamma_{12}$?}

\Stu{The curve $\Sigma_1 \cap \Sigma_2$ from the Triples tab. And it is not defined anywhere on this tab.}

\Kle{Then it is noise for anyone who entered here. Nine-point grey type inside an SVG, no tooltip, referring to an object two tabs back. The rest of the page is scrupulous about naming its terms; this one slipped.}

\Stu{Conceded. Everything else on this tab that needs a definition has a dotted underline and a tooltip. That label has neither because it is drawn text, not HTML.}

\Kle{Also note what the caption gets right, which is not nothing: the two Heisenberg covers share $K_2$ the way the two R\'edei fields share $8081$, and that shared object is why the join is 32 sheets and not 64. That is the sentence that makes the picture worth drawing.}

\walkfig{q8-fieldtower}{The degree-64 field tower, opposite the covers.}

\Stu{The mirror. Three quadratics at the bottom, two biquadratics, the two $D_4$ R\'edei octics $k_{5,8081}$ and $k_{101,8081}$, their compositum at degree 32, and $K$ on top at degree 64 with $\mathrm{Gal}(K/\mathbb{Q}) \cong N_4(\mathbb{F}_2)$. Only invariant-carrying fields are drawn --- degree 64 has vastly more subfields than fit on a page, and the note says so.}

\Kle{Justify the 32. Two octics should give me 64.}

\Stu{The note gives the argument: the two octics intersect exactly in $\mathbb{Q}(\sqrt{8081})$ --- their other quadratic subfields differ, and a $D_4$ octic has no cyclic quartic subfield --- so $8 \cdot 8 / 2 = 32$.}

\Kle{That formula needs both fields Galois over $\mathbb{Q}$, which they are, being $D_4$. The note leaves that step silent. It is true, but a student reading it will think degrees always multiply and divide like that.}

\Stu{Fair. The compositum degree formula is used without its hypothesis being named.}

\Kle{Second thing. Your left branch is labelled ``deg 2, $\sqrt{\alpha}$.'' The right branch, $k_{101,8081}$ over $\mathbb{Q}(\sqrt{101},\sqrt{8081})$, has no label at all --- and its generator is a different R\'edei element, $1009 + 100\sqrt{101}$, which you computed in step 2 and then did not put on the diagram.}

\Stu{That is a genuine omission. The number is on the screen, ten inches above, and the diagram does not connect them.}

\Kle{And the top edge, adjoining $\sqrt\theta$ --- that is the only edge in either tower that carries the new invariant. It is labelled the same size as everything else.}

\begin{knote}
Mirror is honest: base has three components against three ramified primes; $K_4$ and $449$ are both the read object, not the covered one. $\rho(K_4) = $ corner of $N_4(\mathbb{F}_2)$ via the longitude $[[x_1,x_2],x_3]$ --- that is the real content. Counting: $32 \times 2 = 64$ both sides. Gaps: $\gamma_{12}$ undefined on this tab; right branch of field tower unlabelled; compositum-degree hypothesis unstated.
\end{knote}

\walkfig{q9-thread}{The zeta-Delta thread box closing the Quadruples tab.}

\Stu{The thread box, tagged with the zeta-Delta badge that recurs across tabs. It says: $449$ splits completely in the degree-32 compositum --- its six pairwise and four triple symbols, all computed above, are $+1$ --- but the quadruple $-1$ puts $\mathrm{Frob}_{449}$ at order 2 in $\mathrm{Gal}(K/\mathbb{Q}) \cong N_4(\mathbb{F}_2)$: the 64 points fall into 32 orbits of period 2, local factor $(1 - 449^{-2s})^{-32}$. The Pairs box's inert bookkeeping, two levels up.}

\Kle{This is the sentence the whole tab exists to produce, and it is the last thing on the page.}

\Stu{It is placed to be the payoff. And it is state-tracked --- the box is regenerated from the same three branches as the banner, so if the symbol were undefined it would say that no Frobenius order, orbit count, or Euler factor is asserted at $449$, rather than printing a factor anyway.}

\Kle{``The polynomial recording it is the next Alexander-type invariant in the tower.'' Is that computed anywhere?}

\Stu{No. That is a forward reference to the Zeta tab, and it is a claim about what such a polynomial would be, not a computed object. The Zeta tab computes the local factor and stops there.}

\Kle{Then it is a promise. Let us go see whether you keep it.}

\subsection*{Zeta}

\walkfig{z1-intro}{The Zeta tab opening: the packaging claim, with the ``which zeta is which'' note beneath it.}

\Stu{The thesis: everything the previous tabs computed is packaged by a zeta function on each side, and the packaging is the same --- a sum, resummed as a product, in two different ways. Primes play the role of periodic orbits; characters diagonalize Frobenius the way eigenvalues diagonalize the monodromy.}

\Kle{``Periodic orbits of $\mathrm{Spec}(\mathcal{O}_K)$.'' Orbits of what map? There is no dynamical system on $\mathrm{Spec}(\mathcal{O}_K)$.}

\Stu{There is not, and the second paragraph exists to say so. It is headed ``which zeta is which,'' and it states that three distinct objects appear on this page --- the Lefschetz zeta of a monodromy, the Dedekind zeta of a number field, and the Alexander and Iwasawa polynomials --- and that they are related by analogy, not equality. The route in the first paragraph goes through Weil's zeta of varieties over finite fields, where Frobenius really is a map and its fixed points really are the points over $\mathbb{F}_{q^n}$; Artin--Mazur imitated that shape when they named the dynamical zeta. The analogy is on the shape of the resummation.}

\Kle{Then that second paragraph is the most important text on this page, and you have set it in the smallest, greyest type on the page, underneath a bolded sentence that overclaims.}

\Stu{That is right and I do not have a defence. The disclaimer is styled as an aside; the claim it qualifies is styled as a headline.}

\Kle{Invert them. The honest sentence should be the one in bold. What you actually share --- one sum, two product resummations --- is a real and non-obvious observation; it does not need the dressing.}

\walkfig{z2-faces}{The three-faces framework: iterates, eigenvalues, orbits against ideals, characters, primes.}

\Stu{The two columns, three faces each. On the left: sum over iterates, product over eigenvalues of $h_*$ on $H_1(F;\mathbb{Q})$, product over periodic orbits. On the right: sum over ideals, Euler product over primes, product over characters. And between the first two left-hand faces, the resummation engine, shown once: $-\log(1-x) = \sum x^n/n$, so one eigenvalue gives $(1 - \lambda t)^{-1}$. Every sum-into-product identity on the page, both columns, is that one line --- run per eigenvalue, per orbit, or per prime.}

\Kle{The orbit face has a factor $(1-t)^2$ hanging off the left. Explain it, because on a once-punctured torus $H_2$ vanishes and I do not see where a square comes from.}

\Stu{The tooltip says exactly that, and it is the honest scoping on this page. The orbit product counts every fixed point, so it computes the Artin--Mazur zeta, which is a different function from the $H_1$-zeta of the two faces above. With the closed-torus bookkeeping the alternating Lefschetz number is $1 - \mathrm{tr}(A^n) + 1$, every fixed point of a positive-eigenvalue Anosov has index $-1$, so $\#\mathrm{Fix}(h^n) = \mathrm{tr}(A^n) - 2$ --- and the two constants are the $H_0$ and $H_2$ traces. On the punctured fiber $H_2(F) = 0$, so the $(1-t)^2$ would be wrong; the tooltip says the scoping is to the closed torus.}

\Kle{And ``both eigenvalues positive'' is doing work in that hypothesis, not decoration.}

\Stu{The tooltip runs the counterexample: $-A$ is Anosov, but on odd iterates every index is $+1$, so $\mathrm{Fix}(n) = 5, 5, 20, 45, \ldots$ while $\mathrm{tr} - 2 = -5, 5, -20, 45, \ldots$ --- the two sides already differ at $t^1$. Drop positivity and the face is false.}

\Kle{Good. That is a real hypothesis with a real counterexample attached, which is more than most expositions of this manage. Now the layout. Your left column runs the full height of the screen and your right column stops a third of the way down. Half the arithmetic panel is empty. On a page whose entire argument is symmetry, that asymmetry is the first thing the eye reads.}

\Stu{The right column is three faces and a one-line note; the left column carries the convention bridge and the orbit bookkeeping. There is no filler on the right and I did not want to invent any.}

\Kle{Then move something. The Milnor sentence at the bottom left, or the resummation engine, is a shared claim --- it belongs between the columns, not inside one.}

\begin{knote}
Framework is sound. Resummation engine stated once and reused --- correct pedagogy. Orbit face's $(1-t)^2$ properly scoped to the closed torus, with the $-A$ counterexample showing positivity is a real hypothesis. $H_1$-only convention flagged, bridge to alternating-product zeta given. Layout: right column half empty, and the symmetry claim is what the page is selling.
\end{knote}

\walkfig{z3-capstone}{The capstone thread box: the trail of Euler factors computed across the whole app.}

\Stu{The capstone. It reads back the trail: the Pairs box's split and inert local factors, the Triples box's $(1 - 937^{-2s})^{-4}$ at $937$, the Ladder's order-into-orbits reading, the Quadruples box's $(1 - 449^{-2s})^{-32}$ at $449$. Those two factors are recomputed here, not copied --- the Frobenius order at $937$ and the quadruple symbol at $449$ are both rerun to build this sentence.}

\Kle{So if I sabotaged the quadruple symbol, this box would change.}

\Stu{It would say ``$449$ factor --- undefined this run, not asserted.'' The function takes a null and refuses.}

\Kle{``$\Delta_L$ and $f_S$ are, up to units, these zetas.'' What is $f_S$?}

\Stu{The Iwasawa polynomial attached to a finite set $S$ of ramified primes, with finite deck group --- it is defined on the Dictionary tab, with a why-box on the asymmetry.}

\Kle{Which is four tabs away, and this page never says so. Worse: further down this same page you use $f_\Lambda$ for the characteristic power series, and you flag there that $f_\Lambda$ must be kept apart from the residue degree $f$ in your local-factor tables. So this page carries three different $f$'s and disambiguates two of them.}

\Stu{Correct. $f_S$ appears here undefined and unlinked, and the disambiguation the page does make is a hundred lines below the place it is needed.}

\walkfig{z4-spine}{The gallery spine: the first of five row heads.}

\Stu{The gallery is five rows, one per level of the story. Two components --- the pairwise dial, $lk \bmod 2$ against the Legendre symbol. Three components --- the Borromean rings against $k_{13,61}$. Four components --- $B_4$ against the degree-64 field. Five components --- the half-built rung, one column computing and the other naming its missing theorem. Then a contrast pair for what knotting, as opposed to linking, adds. Each card starts closed and leads with its lesson, stated across the analogy.}

\Kle{Cards closed by default. So a visitor sees five headings and no mathematics.}

\Stu{The trade was deliberate: opened, the gallery is roughly six screens of dense output, and the spine is the argument. The row heads are meant to be readable as a single claim --- the levels of linking, in order, with both columns at each.}

\Kle{The spine is good. I can read the whole thesis off five lines, which is more than most talks give me. But default-closed means the app's best material is behind a click, and the row head is styled quieter than the tab's own intro paragraphs. Reverse the emphasis.}

\walkfig{z5-hopf}{The worked Hopf-link card: both positions of the two-component dial.}

\Stu{The smallest case, worked. Face one: the Hopf link is fibered, the fiber is an annulus, $H_1 = \mathbb{Z}$, and the monodromy is the identity, so $\det(tI - h_{1*}) = t - 1$. Torres agrees, computed: $(t-1) \cdot \Delta_{\mathrm{Hopf}}(t,t) = (t-1) \cdot 1 = t-1$, degree one, which is $b_1$ of the annulus. Face two is the dial. $lk(K_1,K_2) = 1$ is odd, so in the double cover of $S^3 \setminus K_1$ the loop $K_2$ lifts to one loop winding twice --- factor $(1 - t^2)^{-1}$ --- against $(5/13) = -1$, computed live: $13$ is inert in $\mathbb{Q}(\sqrt5)$, Euler factor $(1 - 13^{-2s})^{-1}$. Flip the dial: $lk$ even gives two loops, $(1-t)^{-2}$, against $(5/29) = +1$, $29$ split, $(1 - 29^{-s})^{-2}$.}

\Kle{The two Legendre symbols are computed and the two lift counts are not.}

\Stu{Correct. The lift counts are covering-space theory applied to a case anyone can see; the symbols come out of Euler's criterion at run time. The card does not claim otherwise, but it also does not mark which is which as sharply as the higher cards do.}

\Kle{Here is what I actually like. The Hopf link is where this analogy is usually introduced and then abandoned as too small to be interesting. You use it to show the counting rule in its simplest possible instance --- deck order equals loop count times winding --- and then you say, in the closing line, that everything the higher rows do is already here at $n = 2$. That is the correct pedagogical order and almost nobody does it.}

\Stu{The higher rows only climb to deeper covers. The bookkeeping never changes.}

\walkfig{z6-sqrt5}{The zeta of $\mathbb{Q}(\sqrt5)$: three faces converging live under the slider.}

\Stu{An actual Dedekind zeta, three faces at once. The table gives four Euler factors read off Legendre symbols computed live: $5$ ramified, $11$ split, $13$ inert --- that is the Hopf pair from the Pairs tab sitting inside a zeta function --- and $29$ split. Then the slider sets a cutoff $X$, and three numbers update: the sum over ideals, the Euler product over primes, and $\zeta(2) \cdot L(2,\chi_5)$.}

\walkbeat{He leans in on the three numbers.}

\Kle{At $X = 2000$ you show me $1.16145678$, $1.16160385$ and $1.16131818$. Those are three different numbers. They agree to four figures and then part company.}

\Stu{Because they are three different truncations of the same limit, not three evaluations of the same finite sum. The ideal face sums over ideals of norm at most $X$. The Euler face takes complete local factors --- all powers --- for primes up to $X$, so it includes contributions from norms far beyond $X$. The character face truncates $\zeta$ and $L$ at $n \le X$ independently, and the product of two truncations is not the truncation of the product.}

\Kle{Then show me. Move the slider.}

\walkbeat{The student drags to $X = 6000$; the three figures tighten.}

\Stu{They close, and they close at the rate you would expect. The honest version of the caption would say ``three truncations, converging together,'' which is what it does say --- ``converging together,'' not ``equal.''}

\Kle{Then the caption is right and the display is misleading, because eight decimal places invite me to compare eight decimal places. Print four and a residual, or print the truncation each face is using next to it.}

\Stu{Agreed. The precision is beyond what the truncations support.}

\Kle{Why $s = 2$?}

\Stu{Convergence. At $s = 1$ there is a pole; at $s = 2$ all three faces converge fast enough to move visibly under a slider inside a browser.}

\begin{knote}
$\mathbb{Q}(\sqrt5)$ card is the best thing on the tab: three genuinely different computations of one number, all live, and the Legendre symbols driving the Euler factors are the same ones from the Pairs tab. Slider makes convergence visible. But eight printed decimals for three differently-truncated quantities is a display error --- they agree to four and cannot agree to more.
\end{knote}

\walkfig{z7-borro}{The Borromean link-zeta card: fibered, Torres, unipotent, linking at $t = 1$.}

\Stu{The three-component topology column. The Borromean rings are fibered --- Stallings 1978, closures of homogeneous braids, and these close $(\sigma_1\sigma_2^{-1})^3$ --- with fiber the Bennequin surface, so $\chi = 3 - 6 = -3$ and $b_1 = 4$. Torres, computed: collapsing the Triples grid's variables to one and summing by total degree gives $(t-1)^3$, times $(t-1)$ gives $(t-1)^4 = t^4 - 4t^3 + 6t^2 - 4t + 1$, coefficients $1, -4, 6, -4, 1$. Milnor then says the one-variable Alexander polynomial is $\det(tI - h_{1*})$, so every eigenvalue is $1$: unipotent, entropy zero, Lefschetz numbers flat at $4, 4, 4, 4$.}

\Kle{So the punchline is: a link of three unknots has a zeta with no spectral content, everything in the Taylor coefficients at $t=1$. Is that a theorem?}

\Stu{No, and the card says so under a heading that reads ``Scope: exhibited, not a theorem.'' Unknotted components alone do not force it. The counterexample is printed: the $(2,4)$ torus link is two unknots, fibered, $lk = 2$, with $\Delta \doteq (t-1)(t^2+1)$ --- eigenvalues $\pm i$, which are not $1$. They sit on the unit circle so entropy is still zero, but there is spectral content.}

\Kle{Good. So ``linking'' in your slogan means the vanishing-$\bar\mu$ story, not ``components are unknots.''}

\Stu{The card says exactly that in its last clause.}

\Kle{One thing bothers me here. You write $\chi$ for Euler characteristic in the second line and $\chi$ for a Dirichlet character in the column opposite, in the same field of view.}

\Stu{The card interrupts itself to say so --- ``$\chi = $ Euler characteristic here, not this tab's Dirichlet character.'' It is a parenthesis in the middle of a computation, which is the only place it could go, but it is also the third such collision the tab has to defuse: $\gamma$ for orbits against the Triples curve, $\lambda$ for dilatation against the Iwasawa invariant, and now $\chi$.}

\Kle{Three notation collisions, all flagged, none resolved. Flagging is better than nothing. Renaming would be better than flagging.}

\walkfig{z8-borroworked}{The worked Borromean card: two distinct zetas of one link, kept apart.}

\Stu{This card exists because the previous one is easy to misread. Object one is the monodromy zeta on the fiber's homology: traces $4, 4, 4, 4$, and both faces multiplied out to degree eight --- $\exp(\sum 4t^n/n)$ gives $1, 4, 10, 20, 35, 56, 84, 120, 165$ and $(1-t)^{-4}$ gives the same, checked coefficient by coefficient. The right side is computed as $\binom{n+3}{3} = (n+1)(n+2)(n+3)/6$ by literal products --- polynomial growth, no spectral growth. Object two is the cover-orbit factor: in the 8-sheet mod-2 Heisenberg cover of $S^3 \setminus (K_1 \cup K_2)$, $\rho(K_3) = (0,0,1)$ has order 2 by powering, so $K_3$'s preimage is four loops each winding twice --- $(1 - t^2)^{-4}$, the exact twin of $937$'s $(1 - 937^{-2s})^{-4}$.}

\Kle{Two zetas of one link, living in different places. Which one is the analogue of the Dedekind zeta?}

\Stu{The second. The card's lesson line is explicit about it: it is the cover-orbit factor, not the fiber zeta, that is literally the same bookkeeping as $937$'s Euler factor. Deck order equals count times size on both sides.}

\Kle{Then the fiber zeta is the odd one out, and it is the one you compute in most detail.}

\Stu{Because it is the one that has a spectrum to talk about, and the contrast pair at the bottom of the gallery needs it. But you are right that a reader could come away thinking $\zeta_h$ is the mirror of $\zeta_K$. The closing note says ``keep them apart, different objects,'' and I would rather that were the heading than the footer.}

\Kle{The trace stand-in --- ``a stand-in unipotent with the same characteristic polynomial.'' You are not using the actual monodromy matrix.}

\Stu{No. Every quantity used depends only on the eigenvalues, and the card says so each time it appears. The characteristic polynomial is $(t-1)^4$ from Torres and Milnor, both computed; the matrix realizing it is chosen, not derived. That is an honest shortcut but it is a shortcut.}

\walkfig{z9-k1361}{The worked zeta of $k_{13,61}$: eleven local factors, the partial Euler product, the ideal-sum face, and the Artin note.}

\Stu{The arithmetic mirror of the Borromean row, and the densest card on the page. Eleven rows, one per prime: $2, 3, 5, 7, 11, 17, 19, 23, 29, 107, 937$. Each row shows the pairwise level computed --- $(13/p)$ and $(61/p)$ --- then the corner entry, then $f = \mathrm{ord}\,\rho(\mathrm{Frob}_p)$, then the orbit picture, then the local factor. The rule at every unramified $p$ --- which is every $p$ except $13$ and $61$, with $2$ and infinity included --- is $e = 1$, so $fg = 8$: the eight embeddings fall into $g = 8/f$ orbits of size $f$.}

\Kle{$107$ splits completely and $937$ has four orbits of size two. That second one is your triple symbol.}

\Stu{That is the point of the row. $937$'s pairwise data are both $+1$, exactly like $107$'s --- the two look identical at the pairwise level --- and they separate only at the corner: $\alpha$ is a square mod $107$ and a non-square mod $937$. One triple symbol, and the whole Euler product feels it.}

\Kle{Multiplied out, you get $1.067395$ over $44$ primes below $200$. And then a second face agreeing with it. How independent is that agreement?}

\Stu{Not very, and the tooltip concedes it. The ideal sum runs over $200$-smooth norms coprime to $13 \cdot 61$ up to $120{,}000$ --- $151$ terms --- and restricting to $200$-smooth norms makes it converge to exactly the same truncation as the product face. The tooltip says: the two faces then must agree, and the display checks that they do. So it is a cross-check between two code paths, not independent numerical evidence.}

\Kle{Then it is a unit test wearing a mathematical hat, and the admission is inside a tooltip on the word ``smooth.'' Put it in the body text. The agreement to $1.5 \times 10^{-7}$ looks like a discovery and is not one.}

\Stu{Accepted. It does verify something real --- that the ideal-counting recipe $a(p^k) = \binom{k/f + g - 1}{g-1}$ and the Euler factors are the same combinatorics --- but that is a weaker claim than the display implies.}

\Kle{And the third face?}

\Stu{Missing, honestly. $\mathrm{Gal}(k_{13,61}/\mathbb{Q}) \cong D_4$ is nonabelian, so the character face upgrades from Dirichlet $L$-functions to Artin $L$-functions --- stated, citing Artin 1923 and 1930, not computed. The two-component cards run the abelian version because there the group is $\mathbb{Z}/2$ and the characters are on the page.}

\Kle{That is the right call and the right place to stop. Now: this table is unreadable. Your $p = 2$ row wraps its pairwise cell into five lines to say ``$13 \equiv 5$, $61 \equiv 5 \pmod 8$, both inert.'' Six columns in that width was one too many.}

\Stu{The orbit-picture column is redundant with $f$ and the local factor --- it is there for readers who do not want to decode $8/f$. Folding it into the factor column would buy the width back.}

\begin{knote}
$k_{13,61}$ card: $107$ against $937$ is the single best example in the app --- identical pairwise data, separated only by the corner, and the Euler product changes. Partial product $1.067395$. Ideal-sum ``agreement'' is engineered (same truncation by construction) --- admitted only in a tooltip. Artin upgrade for nonabelian $D_4$ correctly stated and correctly not attempted. Table too narrow for six columns.
\end{knote}

\walkfig{z10-b4}{The worked $B_4$ card, at the Milnor-model level.}

\Stu{The four-component topology column, and it is short because the Quadruples tab did the work. $\bar\mu(1234) = 1$, the coefficient of $X_1X_2X_3$ in the Magnus expansion of $\ell_4$, recomputed here. It places $\rho(K_4)$ at the corner of $N_4(\mathbb{F}_2)$ with all $\chi$ and $r$ entries zero. Powered live: order 2. So in the 64-sheet cover $K_4$'s preimage is 32 loops of double length, local factor $(1 - t^2)^{-32}$, the exact twin of the $449$ card's $(1 - 449^{-2s})^{-32}$.}

\Kle{The lesson line claims the quadruple symbol is invisible to every pair and triple, yet the cover hears it. Is that a computation or a slogan?}

\Stu{Both halves are computed, on different tabs. That every pair and triple is trivial is the Quadruples tab's steps 1 and 2b for the arithmetic side and the vanishing lower Magnus coefficients for the topological side. That the cover hears it is the order-2 corner element, powered here. What is not computed is the general statement --- that no pair or triple invariant can ever detect a corner entry. That is the structure of the unipotent group, and the Ladder tab argues it.}

\Kle{So this card is a rendering of results from two other tabs.}

\Stu{Yes, and that is the design of the gallery: each card recomputes rather than caches, but the mathematics is the mathematics of its home tab. This card's own contribution is the twinning --- putting $(1 - t^2)^{-32}$ and $(1 - 449^{-2s})^{-32}$ in one sentence, with the arithmetic side recomputed so that the word ``twin'' is checked and not asserted. If $449$ came out split, the card would say so and would not claim a twin.}

\Kle{Then say that in the card. ``The exact twin'' is doing a lot of work and the reader cannot see that it is conditional.}

\walkfig{z11-449}{The worked degree-64 card at $449$: the order rule brute-forced, and the $\chi$-case table.}

\Stu{The heaviest card. First the order rule of $N_4(\mathbb{F}_2)$, stated: write $M = I + N$, and in characteristic two $(I+N)^2 = I + N^2$ with $N^4 = 0$, so $N^2$ has exactly three entries that can survive --- $\chi_1\chi_2$, $\chi_2\chi_3$, and $\chi_1 r_{234} \oplus r_{123}\chi_3$. Order is 4 if $N^2 \neq 0$, else 2 if $N \neq 0$, else 1. Then it is brute-forced: every one of the 64 elements is literally powered and compared against the rule. Sixty-four out of sixty-four.}

\Kle{You verified a rule about a group of order 64 by checking all 64 elements. That is not a proof technique I would normally applaud, but here it is exactly right --- the group is small enough that exhaustion is a proof.}

\Stu{And it runs at page load, so an edit to the matrix code breaks the display rather than silently changing the answer.}

\Kle{The $\chi$-case table --- eight rows. Some determine $f$ and some do not.}

\Stu{That is the honest part. Each row's verdict is the computed set of orders over the eight fillings of $(r_{123}, r_{234}, q)$. Where $f$ is not determined the row says exactly which deeper entry it needs: $(1,0,0)$ needs $r_{234}$, $(0,0,1)$ needs $r_{123}$, $(1,0,1)$ needs $r_{234} \oplus r_{123}$, and $(0,0,0)$ needs the $r$'s and the corner. Where the $\chi$-products already force $N^2 \neq 0$, the row says order 4 regardless.}

\Kle{So the arithmetic content of the tab lives in exactly one row of that table --- the all-zeros row --- and that is the row where the answer needs the corner.}

\Stu{Which is $449$: $\chi = (0,0,0)$ computed, $r_{123} = r_{234} = 0$ computed, so every entry of $N$ vanishes except possibly the corner, and the quadruple symbol supplies it. Order 2, 32 orbits of size 2, $(1 - 449^{-2s})^{-32}$.}

\Kle{Why no Euler product on this card, when you gave me one for $k_{13,61}$?}

\Stu{Because for a general prime the $r$-entries are its splitting data in the two R\'edei octics, and the app cannot compute those. The closing note says exactly that: only the structure --- the rule, the eight cases --- and the one worked prime are shown, with the general local rule stated for the record.}

\Kle{That is the right refusal. You could have faked a product by pretending the $r$'s were zero everywhere and nobody would have checked.}

\walkfig{z12-five}{The five-component card: $\bar\mu(12345)$ and the 1024-sheet cover, both computed.}

\Stu{The next rung's topology column, running today. $\ell_5 = [[[x_1,x_2],x_3],x_4]$ --- Milnor's model, one commutator deeper --- pushed through the same Magnus machinery already on the page at truncation degree five. Coefficient of $X_1X_2X_3X_4$ is $1$, lower coefficients $0$ and $0$, so $\bar\mu(12345) = 1$ at the model level. Then the counting space: $\rho(K_5)$ at the corner of $N_5(\mathbb{F}_2)$, order $2^{10} = 1024$, and powering the $5 \times 5$ matrix gives order 2. So $512$ loops of double length, factor $(1 - t^2)^{-512}$.}

\Kle{Truncation degree five, and you need a degree-four coefficient. That is a margin of one.}

\Stu{It is, and it is why the card names the truncation rather than hiding it. The machinery was written for the four-component case and the five-component word fits inside it with one degree to spare. Going one rung further would need the truncation raised, and I have not tested that it stays exact.}

\Kle{Say it plainly: this is the last rung this code can reach.}

\Stu{With the current truncation, yes.}

\Kle{And what have you actually demonstrated? That a commutator you wrote down has the coefficient you expected, and that a corner element of a matrix group has order two. Neither is surprising.}

\Stu{Neither is surprising, and the card's claim is correspondingly modest --- the lesson line says the ladder does not stop, and that what is missing on the other side is a theorem, not a computation. The value is not the number; it is that the topological column can be produced at all at a level where the arithmetic column cannot.}

\walkfig{z13-fiveblank}{The deliberately blank arithmetic cell of the five-component row.}

\Stu{And this is the other half of that row. No quintuple symbol exists to compute. The analogue of Amano's independence theorem for $[p_1, \ldots, p_5]$ --- a normalization making the bracket a function of the five primes alone, reading splitting in a degree-1024 field with Galois group $N_5(\mathbb{F}_2)$ --- is not in the literature the app cites. Until that theorem exists the cell stays blank rather than pretend.}

\Kle{``Not in the literature this app cites'' is a careful phrase. It is not the same as ``does not exist.''}

\Stu{It is deliberately weaker, because I cannot certify the second. I can certify what I read.}

\Kle{Good. Then the blank is honest twice over --- about the mathematics and about the limits of your search. ``The app's one citation to an absence --- the citation style one hopes goes stale'' is a nice line, and unlike most nice lines on this page it is earned. Leave the cell empty. Do not let anyone talk you into filling it with a diagram of what the theorem would look like.}

\Stu{That was the temptation and I resisted it.}

\begin{knote}
Five-component row is the best-judged thing in the app. Topology computed ($\bar\mu(12345) = 1$, order 2, $512$ double loops); arithmetic cell blank with a stated reason. Margin: Magnus truncation is degree five and the needed coefficient is degree four --- one to spare, so this is the ceiling of the current code. Blank is scoped to ``literature this app cites,'' not to ``does not exist.'' Correct.
\end{knote}

\walkfig{z14-fig8}{The figure-eight contrast card: knotting moves the spectrum.}

\Stu{The contrast pair. $4_1$ is fibered like the Borromean rings, but knotted: the monodromy acts on $H_1$ by the Anosov matrix $A = [[2,1],[1,1]]$ with dilatation $(3+\sqrt5)/2 = \varphi^2$. The table gives $\mathrm{tr}(A^n) = 3, 7, 18$, then $\mathrm{Fix}(h^n) = \mathrm{tr} - 2 = 1, 5, 16$, then primitive orbit counts $1, 2, 5$. Below it, three identities checked coefficient by coefficient to degree eight: the two faces of $\zeta_h$ agree at $1, 3, 8, 21, 55, 144, 377, 987, 2584$ --- alternate Fibonacci numbers --- the two faces of the Artin--Mazur zeta agree, and $(1-t)^2$ times the $H_1$ face equals the orbit face.}

\Kle{The third identity is the one that matters. It is the $(1-t)^2$ you asserted in the framework, now verified.}

\Stu{Verified nine coefficients deep. And the orbit-face check does something extra: it verifies integrality of the M\"obius counts. Whole-number orbit counts are the statement that iterates organize into closed orbits --- if the identity held numerically but the orbit counts came out fractional, the interpretation would be empty.}

\Kle{That is a real point and it is buried in a grey note. Now look at your own display. You have printed ``$\prod\_d (1 - t^{\wedge}d)^{\wedge}\{-o\_d\}$'' with literal braces and carets, next to a line that renders $\exp(\sum \mathrm{Fix}(h^n) \cdot t^n/n)$ properly typeset. And three inches above, in the introduction to this card, you have written zeta-sub-h as ``zeta underscore h'' twice.}

\Stu{Those are two hard-coded label strings that never got converted to markup, while everything around them did. It is a real defect and it is at the visual center of the card that carries the tab's punchline.}

\Kle{It undercuts you more than you think. A reader who cannot tell whether your notation is deliberate will not trust your arithmetic either.}

\Stu{The honesty note below is also load-bearing and I should say it: the fiber is a once-punctured torus, but the fixed-point counts are taken for the same matrix on the closed torus, where $\mathrm{Fix} = |\det(A^n - I)| = \mathrm{tr}(A^n) - 2$ and the orbit bookkeeping is clean. The $H_1$-action, which is all the faces use, is identical.}

\Kle{And the dichotomy line at the bottom: entropy $\log \varphi^2 \approx 0.962$ against the Borromean box's flat $4, 4, 4$. That is the whole tab in one comparison. It should not be the last box on the card.}

\walkfig{z15-knotted}{The knotted-primes card: regular against irregular, computed at load.}

\Stu{The arithmetic mirror of knotting. Mazur's dictionary: a knot's Alexander polynomial against a prime's Iwasawa polynomial for its own cyclotomic $\mathbb{Z}_p$-tower, with $\lambda$ its degree. Regular prime, trivial polynomial, unknotted --- Kummer's criterion plus Iwasawa's theorem. Irregular, knotted. The test is Bernoulli, run mod $p$, and the cast is judged at load in four milliseconds: $5, 13, 29, 61, 449, 937$ all regular, and $101$ irregular --- $B_{68} \equiv 0 \pmod{101}$.}

\Kle{Running the Bernoulli recurrence mod $p$ needs the denominators to be prime to $p$.}

\Stu{Von Staudt--Clausen, and the card cites it with the argument sketched: $p$ divides the denominator of $B_k$ only if $(p-1) \mid k$, which is impossible for even $k \le p - 3$. So those $B_k$ are $p$-integral, the recurrence reduces, and $B_k \equiv 0$ is equivalent to $p$ dividing the numerator.}

\Kle{And $101$ is one of Amano's four primes.}

\Stu{That is the card's jewel, and the sentence I am most careful about: the app contained a knotted prime all along and the linking story never noticed. These are independent dials --- the $\ell = 2$ linking tower of $\{5, 8081, 101, 449\}$ and the $p = 101$ cyclotomic tower are different covers of different things. The coincidence is that $101$ wears both hats, not that one caused the other.}

\Kle{Do not oversell it and you have a genuinely charming observation. Oversell it and it becomes numerology. The sentence as written is on the right side of that line.}

\walkbeat{He scans the cast table.}

\Kle{Where is $8081$? It is one of the four and it is not in the table.}

\Stu{Behind that button. Deciding $8081$ needs the recurrence run to $B_{8078}$, which is too slow for page load, so it is chunked and asynchronous with a progress tick. The other seven are cheap enough to judge at load.}

\Kle{So three of Amano's four primes are judged in front of me and the fourth requires that I notice a small grey button below a table. That is the pattern I keep hitting on this app: the supporting evidence is one click away and the click is invisible.}

\Stu{That is the fairest single criticism of the whole design.}

\Kle{And the $\lambda$ collision --- you use $\lambda$ for the dilatation on the card opposite.}

\Stu{Flagged in the first sentence: ``not the figure-eight card's dilatation $\lambda$ opposite: same letter, different object.'' Flagged, again, rather than fixed.}

\walkfig{z16-presmat}{The two presentation matrices: $tI - A$ for the Alexander module, and the $1 \times 1$ $\Lambda$-side.}

\Stu{The closing pair. On the left, the module itself: for a fibered knot the Alexander module is presented by $tI - h_{1*}$, by Milnor's theorem that the infinite cyclic cover is $F \times \mathbb{R}$. For $4_1$ that is the two-by-two matrix with entries $t-2$, $-1$, $-1$, $t-1$, and the determinant is computed live by polynomial arithmetic: $t^2 - 3t + 1 = \Delta_{4_1}$, the polynomial the figure-eight card verified as the zeta.}

\Kle{And the right.}

\Stu{The mirror module at the level the literature pins. The Iwasawa module of an irregular prime is a $\Lambda = \mathbb{Z}_p[[T]]$-module; for $p = 37$ the relevant $\omega$-eigenspace has $\lambda = 1$, $\mu = 0$ --- $\mu = 0$ by Ferrero--Washington, the $\lambda$-value from the standard tables --- so the characteristic polynomial is linear and the smallest faithful presentation is one-by-one: $\Lambda / (f_{37}(T))$.}

\Kle{You have drawn me a one-by-one matrix containing a symbol whose value you do not know.}

\Stu{And the note says exactly that: what the cited results pin are the matrix's size and its determinant ideal --- the characteristic ideal is $(L_p)$ by the Main Conjecture. The entries need class-group tower data the app does not carry. It is displayed schematically and the display says so, with the contrast drawn explicitly: the two-by-two twin at left is fully computed.}

\Kle{Then the two halves of this figure are not the same kind of object. One is a computation and one is a diagram of a computation somebody else did. The layout presents them as a matched pair.}

\Stu{The visual symmetry outruns the epistemic symmetry. That is true here and it is true of the whole app --- the two-column format asserts a parity between the sides that the individual boxes then have to walk back one at a time.}

\Kle{That is the most useful thing you have said in an hour. The format is an argument, and it is an argument you have to keep apologizing for. On $37$ specifically --- why not $101$, which you just told me is irregular and is one of your four?}

\Stu{Because $37$ is the classical first case with tabulated invariants I can cite. For $101$ I would need the $\lambda$-value and I did not verify one I trust.}

\Kle{Then say ``$37$, because its invariants are tabulated and $101$'s are not to hand.'' That sentence costs you nothing and buys the reader the reason the two cards do not line up.}

\begin{knote}
Closing pair: left is fully computed ($tI - A$, $\det = t^2 - 3t + 1 = \Delta_{4_1}$, live). Right is schematic --- size and determinant ideal cited (Ferrero--Washington $\mu = 0$, $\lambda = 1$ from tables), entries unknown, admitted. $p = 37$ used, not $101$, and the reason is not given.

Overall, Zeta tab: the sum-two-resummations claim is real and is carried through every card. Scope discipline is unusually good --- the $(1-t)^2$ hypothesis, the $(2,4)$ torus counterexample, the Artin upgrade left undone, the blank fifth cell. Weaknesses are presentational and they are systematic: the disclaimer is smaller than the claim, the raw-TeX labels sit at the center of the punchline card, and the evidence is repeatedly one invisible click away. The two-column layout asserts a symmetry the content spends its time qualifying.
\end{knote}\subsection*{Q and A}

\walkfig{a1-intro}{The Q and A tab opens on a single promise about vocabulary.}

\Stu{This tab is the skeptic's index. Fourteen questions, all closed at load. The rule in that intro line is the rule the whole app runs on: every answer either jumps into a tab where the claim can be recomputed on the spot, or it carries a name and a year --- and the three words \emph{theorem}, \emph{computed live}, and \emph{analogy} never trade places.}

\Kle{Fourteen questions a skeptical manifold theorist asks. Written by you, in my voice, before I got here.}

\Stu{Yes. That is the obvious objection and I will not dress it up: I chose the questions, so I chose the ones I could answer. What I can defend is that none of the fourteen has a soft answer inside it. Three of them are failure lists.}

\Kle{Fine. I will supply my own questions as we go and we will see which ones you anticipated. Open the first.}

\walkfig{a2-q1}{Question one: why should a prime resemble a knot at all.}

\Stu{Two theorems and one leap, kept strictly apart. Theorem: $\mathrm{Gal}(\bar{\mathbb{F}}_p/\mathbb{F}_p) \cong \hat{\mathbb{Z}}$, topologically generated by Frobenius --- the same profinite group as the profinite $\pi_1$ of the circle, so $\mathrm{Spec}\,\mathbb{F}_p$ is a $K(\hat{\mathbb{Z}}, 1)$ exactly as $S^1$ is a $K(\mathbb{Z},1)$. Theorem: $\mathrm{Spec}\,\mathbb{Z}$ behaves three-dimensionally --- etale cohomological dimension 3, with Artin--Verdier duality playing its Poincare duality, Artin--Verdier 1964 and Mazur 1973. And then the leap, in bold and labelled Analogy: a circle sitting in a 3-space is a knot. Mazur's observation, unpublished notes around 1963 to 1964.}

\Kle{Read me the sentence after the leap.}

\Stu{``Asserts no embedding and no isotopy; that leap \emph{is} the dictionary.''}

\Kle{Good. That is the sentence that matters and I am glad it is on the page rather than in your mouth. Because the honest statement of the situation is that you have a $K(\pi,1)$ on one side, a duality of the right formal shape on the other, and nothing whatsoever connecting them --- no map, no embedding of $\mathrm{Spec}\,\mathbb{F}_p$ into $\mathrm{Spec}\,\mathbb{Z}$ that behaves like an embedding of $S^1$ into $S^3$. There is a map of schemes, and it is not an embedding in any sense I would use the word.}

\Stu{Agreed, and the app never claims one. What survives the leap is that both duality statements are top-degree 3, so the pairings have matching arities --- which is what lets a linking number be defined on both sides at all.}

\Kle{That is the right defence. Say it that way. ``Matching arities'' is a claim I can check; ``a prime is a knot'' is a slogan.}

\begin{knote}
Q\&A tab. Q1 is honest where it counts: theorem / theorem / analogy, in bold, in that order. He conceded instantly that there is no embedding. The real content is that both dualities are top-degree 3, so the pairings have the same arity --- which is what a linking number needs. Tell him to promote that sentence into the answer; it is currently only in his voice, not on the page.
\end{knote}

\walkfig{a3-q9}{Question nine open: are the two zeta functions the same object?}

\Stu{The answer starts with the word ``No,'' which is the point. The Lefschetz zeta of a monodromy and the Dedekind zeta of a number field are related by analogy, not equality --- and the Zeta tab says so in a note above the two columns, before you can object. What the columns actually share is a mechanism: one sum resummed as a product two ways. Iterates, eigenvalues, orbits on the left; ideals, characters, primes on the right, and every instance reduces to $-\log(1-x) = \sum x^n/n$.}

\Kle{Where does equality ever get asserted, then?}

\Stu{Only within a column, never across. Two up-to-units identities are theorems, each inside one side of the mirror: $\tau_L \doteq \Delta_L$, Milnor 1962 and Turaev 1986, and the Main Conjecture $f \doteq L_p$, Mazur--Wiles 1984, with $\ell = 2$ by Greither 1992. And the one numerical equality on the tab is $\zeta_{\mathbb{Q}(\sqrt{5})}(2)$ computed three ways --- ideal sum, Euler product, character product --- converging to the same decimals, live.}

\Kle{Three ways to compute one number is not three theorems, it is one theorem and two implementations. But I take the point of the display, which is that you are willing to be caught.}

\walkbeat{He reaches over and opens the next question down himself.}

\Kle{This is the one I actually wanted. ``What exactly is computed live, and what am I being asked to take on faith.'' Read me the faith list.}

\Stu{Cited, not computed: fiberedness, Stallings 1978; $\bar\mu(123) = \mathrm{lk}(\gamma, K_3)$, Cochran 1990 and Mellor--Melvin 2003; the model longitudes, Milnor 1957; Redei's ramification statement, 1939; Amano's independence theorem, Tohoku 2014; Milnor--Turaev, Mazur--Wiles, Greither. Computed: every Legendre symbol by Euler's criterion, Tonelli--Shanks square roots, Redei symbols with the normalization enforced, Magnus coefficients, matrix orders by literal powering, Bernoulli numbers mod $p$.}

\Kle{And the flagged gap.}

\Stu{One, on the Quadruples tab: $\bar\mu(1234)$ is computed for Milnor's model $B_4$, not for the drawn diagram. The Wirtinger presentation is not implemented and the honesty note under the picture says so.}

\Kle{Here is my complaint about this entire tab, and I want it recorded. Not one number appears on it. This is the tab whose whole subject is what is computed live, and it computes nothing --- it is fourteen paragraphs of prose with buttons. Every claim's support is one click away and none of it is visible. If I read only this tab I would have to take the word ``computed'' entirely on trust, which is precisely the posture the tab exists to refuse.}

\Stu{That is fair and I have no answer to it. The tab was built as an index, and an index that indexes computation is not computation.}

\Kle{Put one live number on it. One. Recompute $(13/61)$ in the margin of question one and the tab stops being a brochure.}

\walkfig{a4-breakage}{The last question: where the dictionary breaks, five fractures.}

\Stu{Five, each flagged in the app at its point of use. One: the archimedean place --- arithmetic carries $S \cup \{\infty\}$ bookkeeping everywhere, and the whyFS box on the Dictionary tab shows it biting at $\ell = 2$, where the totally real tower drops to $\mathbb{Z}/4 \times \mathbb{Z}/2$; the knot side has no mirror for $\infty$ past the $p \equiv 1 \pmod 4$ orientability convention. Two: no arithmetic hyperbolic geometry --- the figure-eight complement carries a volume and dilatation $\varphi^2$, and no prime carries anything like a volume. Three: the $B_4$ gap. Four: the ring mismatch --- $\Delta_L$ lives over $\mathbb{Z}[t_1^{\pm 1}, \ldots, t_r^{\pm 1}]$ with deck group $\mathbb{Z}^r$, and over $\mathbb{Q}$ no $\mathbb{Z}^r$-cover exists. Five: it is not a functor.}

\Kle{Stop at four. Say the fourth one again, slowly, because it is the only one on the list that is a theorem rather than a shrug.}

\Stu{Over $\mathbb{Q}$ there is no $\mathbb{Z}^r$ deck group unramified outside the $p_i$. The whyFS box computes the $f_S$ deck group finite --- $\mathbb{Z}/4 \times \mathbb{Z}/4$ for $S = \{13, 61\}$ at $\ell = 2$. And the one $\mathbb{Z}_p$-tower $\mathbb{Q}$ does have, the cyclotomic one, ramifies at $p$ itself.}

\Kle{So the two polynomials do not merely happen to live over different rings --- they are \emph{forced} apart, by finiteness on one side. That is a much stronger and much better statement than ``the fit is imperfect,'' and it is buried as item four of five in a paragraph. Fractures one, two, three and five are honest caveats. Four is a theorem about why the analogy cannot be upgraded. Lead with it.}

\Stu{Noted. It is currently ranked by where it appears in the app, not by weight.}

\Kle{And number five --- ``it is not a functor.'' Who is that for? A manifold theorist hears that and asks which category you failed to construct. A number theorist may not hear anything at all. It is the correct statement, but it is the only one of the five that is a claim about mathematics rather than about your app, and it gets one sentence.}

\begin{knote}
Fracture (4) is the real one and it is buried: over $\mathbb{Q}$ no $\mathbb{Z}^r$-cover unramified outside $S$ exists, $f_S$ deck group finite ($\mathbb{Z}/4 \times \mathbb{Z}/4$ for $S = \{13,61\}$, $\ell = 2$), and the genuine $\mathbb{Z}_p$-tower ramifies at $p$. So $\Delta_L$ and the Iwasawa polynomial are forced over different rings --- not a mismatch of taste, a finiteness theorem. This is the sharpest thing on the tab. It is item four of five in a run-on paragraph. Also: this tab computes nothing. Tab about live computation, zero live computation.
\end{knote}

\walkfig{a5-list}{All fourteen questions, closed --- the tab as it loads.}

\Stu{This is the resting state. Everything shut, so the tab reads as a contents page rather than a wall.}

\Kle{Then let me review it as a contents page. They are not numbered. You told me in the first line there are fourteen, so I have to count them myself to find the fourteenth, and there is no way to open them all at once --- if I want to read this tab straight through I click fourteen times. That is a small thing and it is the difference between a document and a toy.}

\Stu{There is no expand-all control. Correct.}

\Kle{Second thing, and this one is not small. Count the buttons that jump out of this tab.}

\Stu{Twenty, across six destinations --- Dictionary, Pairs, Triples, Ladder, Quadruples, Zeta.}

\Kle{Six. There are seven other tabs. Not one of these fourteen answers sends me to the Experimental tab, which is where you keep the things you are least sure about. Your skeptic's index routes the skeptic away from the frontier. If question thirteen --- the five-component one --- had a button to it, I would have gone there under my own steam instead of because you walked me.}

\Stu{That is a genuine omission. Question thirteen ends by saying Amano's quadruple is the frontier, and the tab that argues about frontiers is the one it does not link.}

\Kle{One more. Question six --- ``everything in this app is mod 2.'' Its answer uses the phrase ``the $f_S$ deck group'' twice and never says what $f_S$ is. I can guess: the maximal pro-$\ell$ extension unramified outside $S$. But I am guessing, and I am the audience.}

\begin{knote}
Contents page: unnumbered, no expand-all, fourteen clicks to read. Twenty jump buttons, six destinations, none to Experimental --- the index avoids the frontier. $f_S$ used twice in Q6, defined nowhere on the tab. Overall: the honesty vocabulary holds up across all fourteen. I could not find a place where \emph{theorem} was doing work that only \emph{analogy} had earned. That is rarer than it sounds.
\end{knote}

\subsection*{Experimental}

\walkfig{x1-banner}{The frontier banner: nothing here is load-bearing.}

\Stu{The banner sets the rules for the whole tab. Nothing here is load-bearing for the other seven. Every claim below is marked one of four ways --- \emph{theorem}, \emph{computed live}, \emph{empirical} meaning machine-checked but not derived, or \emph{proposal / conjectural}. Two threads: what the zeta function counts when you change the generator, and what happens to the dictionary at $\ell = 3$.}

\Kle{The four-way labelling is the right idea and I want to be clear that I approve of it. Most people have three tiers at best. ``Empirical, machine-checked but not derived'' is a category most mathematicians refuse to admit exists in writing, and you have given it a name and a badge.}

\Stu{It exists whether or not it is admitted.}

\Kle{It does. Now --- look at how it is set. That paragraph is the constitution of the tab and it is rendered in small grey centred italics. It is the least legible block on the page and it governs everything below it. Whoever chose that style was thinking ``modest,'' and modesty in typography reads as ``skippable.''}

\Stu{The style class is shared with every honesty box in the app --- point-eight-five rem, muted, italic, centred. It is applied to three blocks on this tab and two of them are the longest paragraphs here.}

\Kle{Then it is a systematic error, not a slip. Your most important prose is your least readable prose.}

\walkfig{x2-generators}{The two ramified generators, side by side: the trefoil worked live against the Iwasawa twin.}

\Stu{Left column, the knot story. Feed the Lefschetz formula the monodromy $\varphi_*$ of a fibered knot acting on $H_1$ of its fiber --- a local, ramified generator. Milnor 1968: you get the Alexander polynomial. The trefoil is fibered over $S^1$ with fiber a once-punctured torus and monodromy the $SL_2(\mathbb{Z})$ matrix $[[1,1],[-1,0]]$, characteristic polynomial $\Delta(t) = t^2 - t + 1$. That polynomial is cited. Everything in the table is computed from it here: $\#H_1(M_n) = \prod_{\zeta^n = 1}|\Delta(\zeta)|$, giving 3, 4, 3, 1 at $n = 2, 3, 4, 5$, and infinite at $n = 6$.}

\Kle{Three, four, three, one, infinite.}

\walkbeat{He stops the walkthrough and looks at the column of integers for several seconds.}

\Kle{Those are the Brieskorn spheres. $\Sigma(2,3,n)$. Your $n = 2$ row is the lens space $L(3,1)$, first homology $\mathbb{Z}/3$. Your $n = 3$ is $\Sigma(2,3,3)$ with $\mathbb{Z}/2 \oplus \mathbb{Z}/2$, order four. And your $n = 5$ row, the one you have printed as the integer 1 --- that is the Poincare homology sphere. $\Sigma(2,3,5)$. The most famous 3-manifold in my subject is sitting in your fourth table row labelled ``1'' and you have not named it.}

\Stu{I had not made that identification. The app computes the order and stops.}

\Kle{And $n = 6$ infinite, because $\Delta = \Phi_6$ vanishes at a primitive sixth root --- $\Sigma(2,3,6)$, which is not even hyperbolic, it is the Nil one. Your table is a complete tour of the trefoil's branched covers and it presents itself as a column of arithmetic.}

\Stu{Which brings up something you should say back to me, because I think you are about to.}

\Kle{I am. Nowhere on this tab do you write the word ``branched.'' You say ``the whole tower of cyclic covers,'' and $M_n$ is never defined. But that product formula is the formula for the $n$-fold cyclic \emph{branched} cover of $S^3$ branched over $K$. The unbranched cyclic covers have different homology --- the infinite cyclic cover of the trefoil has $H_1 = \mathbb{Z}^2$ as a group, not order 3. A reader who takes your word ``cyclic covers'' literally will compute the wrong thing.}

\Stu{You are right, and I have checked: ``branched'' appears exactly once in the entire file, in an unrelated Dictionary tooltip about ramification. On this tab it is missing and $M_n$ is undefined. That is a straightforward defect.}

\Kle{It is the good kind of defect --- the numbers are right and the word is missing. Now the right column.}

\Stu{The same formula with the other ramified generator: $\gamma$, the topological generator of a $\mathbb{Z}_p$-extension acting on the Iwasawa module $H_\infty = \varprojlim \mathrm{Cl}(k_n) \otimes \mathbb{Z}_p$. Iwasawa: you get $I_p(T) = \det(T \cdot \mathrm{id} - (\gamma - 1))$ up to $p^\mu$, with $\#H_n = p^{\lambda n + \mu p^n + \nu}$. And the quantity that matches on the topological side is not the order but $v_p(\#H_1(M_{p^n}))$ --- the third column at left. For the trefoil at $p = 2$ that column is identically zero, because 3 and 3 are odd, so $\lambda = \mu = \nu = 0$. The asymptotic in its degenerate case.}

\Kle{You have matched two things by exhibiting an instance where both are zero. That is not a match, that is a coincidence of vanishing.}

\Stu{Agreed, and I would not call it evidence. What it does is force the right statement: the correspondence has to be with $v_p$, not with the integer, because $\#H_1$ need not be a $p$-power at all --- your $n = 3$ row is 4, and 4 is not a power of 2 in any way that would help if $p = 3$. The degenerate case is what makes the shape of the claim visible. It is not what supports it.}

\Kle{Accept that. Now, this column has one genuinely important sentence in it and it is not visible. Where is the statement that $\gamma$ is \emph{not} analogous to Frobenius?}

\Stu{In a hover tooltip. The words ``not to Frobenius'' are underlined and the explanation appears on hover: $\gamma$ generates the ramified tower itself, Frobenius generates an unramified local extension --- the loop around a puncture disjoint from the branch locus. Conflating them is the commonest error in reading the dictionary.}

\Kle{Your own tooltip calls it the commonest error, and you have put the correction behind a mouse. Also --- look at the shape of this screen. Your left column runs to the bottom of the viewport and your right column stops halfway, so there is a column of empty page next to your most delicate paragraph. It reads as though the right side ran out of things to say.}

\begin{knote}
Best beat of the tab and he did not see it: his trefoil table is $\Sigma(2,3,n)$ for $n = 2,3,4,5,6$ --- $L(3,1)$, $\mathbb{Z}/2^2$, $\mathbb{Z}/3$, the \emph{Poincare sphere}, and Nil. He prints ``1'' where $\Sigma(2,3,5)$ belongs. Name them; the whole point of this app is that the two columns illuminate each other, and here the topological column has famous names he is not spending.

Defect: ``branched'' never appears, $M_n$ never defined; the product formula is for the branched covers. Numbers right, word missing.

$\lambda = \mu = \nu = 0$ is agreement by mutual vanishing. He conceded it unprompted and gave the better reason for it (matching must be with $v_p$, not the order, since $\#H_1$ need not be a $p$-power). Correct.

The load-bearing sentence of the tab --- $\gamma$ is not Frobenius --- is in a hover tooltip.
\end{knote}

\walkfig{x3-pivot}{The pivot box: change the generator, change the subject.}

\Stu{This is the hinge of the tab. Feed the Lefschetz formula the ramified generator, $\varphi_*$ or $\gamma$, and it returns $\Delta_K$ and $I_p$ --- knot-type, depth two, blind to $\bar\mu(123)$. Feed it the unramified generator, Frobenius --- the loop around a puncture disjoint from the branch locus --- and the same formula is an Artin $L$-function, which is where the Redei symbol and every higher-linking invariant in this app already live.}

\Kle{Say the second half of the box.}

\Stu{The app read back against itself: Pairs and Triples pair against Frobenius, so they are the link column, unramified. The Zeta tab's figure-eight monodromy card and the knotted-prime card's Iwasawa polynomial use $\varphi_*$ and $\gamma$, so they are the knot column, ramified. Same formula, two seats --- a distinction the other tabs use without naming.}

\Kle{That is the best paragraph in the app and it is on the tab you told me was not load-bearing.}

\Stu{I know.}

\Kle{No, I want you to hear the size of it. You spent seven tabs computing linking numbers and Iwasawa invariants side by side, and the actual reason they are different kinds of object --- one generator ramified, one not --- is quarantined on the frontier tab under a banner saying ``nothing here is load-bearing.'' It is load-bearing. It is the load. Every reader who leaves after the Zeta tab leaves thinking $\Delta_K$ and $\bar\mu(123)$ are the same flavour of invariant, and they are not: one is depth two and one is not.}

\Stu{The banner is about epistemic status, not about importance --- nothing here is \emph{relied on} by the other tabs.}

\Kle{Then your banner is measuring the wrong axis, and it is costing you your best idea. Move this box to the Dictionary tab and label it theorem-plus-definition, because that is what it is. The conjectural material can stay out here.}

\walkfig{x3b-cracks}{The two honest cracks, stated because they are load-bearing.}

\Stu{Two, and both are conceded rather than defended. First: the indeterminacy does not match. Milnor's torsion $\doteq \Delta_K$ holds up to $\mathbb{Z}[t,t^{-1}]^\times = \{\pm t^n\}$, a tiny group. The Main Conjecture's ``up to units'' is up to $\Lambda^\times$ where $\Lambda = \mathbb{Z}_p[[T]]$ --- every power series with unit constant term. Enormously larger. The triangle repeats; its slack does not; and the note says outright that it does not know a reason it should.}

\Kle{That is the correct thing to be worried about and almost nobody writes it down. $\{\pm t^n\}$ is countable and discrete. $\Lambda^\times$ is a pro-$p$ group of infinite rank. An identity ``up to $\Lambda^\times$'' pins a principal ideal, and an identity up to $\{\pm t^n\}$ pins a polynomial to sign and shift. Those are not the same strength of statement by any measure, and if the analogy were structural rather than formal, the slack ought to correspond too.}

\Stu{That is exactly the sentence the box is trying to say and yours is better.}

\Kle{Second crack.}

\Stu{The depth-three statement is a proposal, not a theorem. The reading offered is that the order of vanishing of the Artin $L$-function of the tower generated by the arithmetic Seifert surface $\sqrt{\alpha}$ is the number-theoretic analogue of Cochran's derived invariant, and that the Redei symbol sits inside such an order of vanishing.}

\Kle{``One rung below the machine-checked depth-four case.'' What is the depth-four case?}

\Stu{The $d$-bit law in the box below.}

\Kle{Then say so. The phrase ``depth-four'' appears twice on this tab and is defined neither time, and the thing it refers to is called something else entirely where it appears --- I have just read the next box and it is called ``the one machine-checked instance'' and ``a governing bit.'' You have three names for one object on one screen. And ``depth two'' you did define, in the left column, so a reader has every reason to think depth-four means four link components. Does it?}

\Stu{No. It does not mean four components.}

\Kle{Then that is a term doing real work with no definition and a misleading neighbour. Fix it or drop it.}

\walkfig{x4-redei}{A second Borromean prime triple, computed live: $(13, 17, 53)$.}

\Stu{This is a live computation, not a citation. Pairwise: $(13/17) = +1$, $(13/53) = +1$, $(17/53) = +1$ --- all three, the arithmetic unlink, and all three primes $\equiv 1 \pmod 4$. Redei witness $(x,y,z) = (-15, 4, 1)$: the identity $x^2 - 13y^2 - 17z^2$ evaluates to 0 on screen, $y$ is even, $x - y \equiv 1 \pmod 4$. Conjugates of $\alpha$ mod 53 come out 45 and 31, both Legendre $-1$, so $[13,17,53] = -1$.}

\Kle{Pairwise unlinked, triply linked, on a triple I have not seen before. Good --- that is the answer to the cherry-picking question your Q and A tab gave in prose, given here in numbers instead. Why is $x$ negative?}

\Stu{Some pairs have no normalized certificate with $x$ positive. The engine accepts any sign on the certificate coordinates; only the primes must be positive.}

\Kle{And it is displayed as ``-15'' with a hyphen while your symbol value four lines later is ``$-1$'' with a proper minus sign. On the same card. I mention it because you have clearly been careful about typography everywhere else, so this one is going to look like the number came from somewhere else --- which, in a document about what is computed live versus what is quoted, is exactly the wrong impression to give.}

\Stu{It is a raw join of the certificate array. It should go through the same formatter as everything else.}

\Kle{Now the sentence underneath: it is $\mathrm{Frob}_{53}$ paired against $\sqrt{\alpha}$ that decides it. Unramified generator, doing link theory. That is the pivot box being cashed out one screen later, and it is the first time on this tab that the abstract distinction has produced a number. Say that when you present this --- the reader should see the pivot pay.}

\walkfig{x5-dbit}{The $d$-bit box: the empirical law, this app's own pair, and the normalization caution.}

\Stu{Three blocks, three different epistemic statuses, and they are labelled. Block one is empirical and is not computed here: the laboratory reports, on 8091 pairs, that a governing bit $d\{p,q\} \in \{1,3\}$ defined by the sign of the Redei witness $\alpha$ --- $d = 1$ when $F_{p,q} = \mathbb{Q}(\sqrt{p}, \sqrt{q}, \sqrt{\alpha})$ is totally real, $d = 3$ when totally imaginary --- matches an a priori independent analytic invariant, the order of vanishing of the Artin $L$-function: $d = 1$ if and only if $\mathrm{ord}_{s=0} L(s, \rho_{p,q}) = 2$, and $d = 3$ if and only if the order is 0. The equality of two independently defined quantities is the claim. No derivation is asserted.}

\Kle{Why one and three? Not zero and one, not plus and minus. What are you indexing?}

\Stu{The box does not say, and I should be honest that I cannot reconstruct it from the app. The labels come from the source note.}

\Kle{Then on this page they are decoration. A bit with two values called 1 and 3 invites me to think there is a $d = 2$ somewhere, and there is not. Same for ``a Stark rank'' in parentheses --- that is three words standing in for a theory, and this tab has already shown it is willing to spend a paragraph on things that matter.}

\Stu{Block two is computed live: our own pair. $(13, 61)$ with the shipped normalized certificate $(23, 6, 1)$ gives $\alpha = 23 + 6\sqrt{13} = 44.633$ and $\bar\alpha = 1.367$. Both positive, so totally positive, so $k_{13,61}$ is totally real and $d\{13,61\} = 1$.}

\Kle{And block three is the one you should lead with.}

\Stu{Block three is a caution the app found while checking, and states rather than hides. The sign is only well defined once Redei's normalization is imposed. For $(13,17)$, every one of the 10 normalized witnesses with $z \le 12$ gives $\alpha$ totally negative --- 10 negative, 0 positive, computed at load --- so $d = 3$. But the illustrative witness $(9,1,2)$ quoted in the source note is primitive and non-degenerate and \emph{not} Redei-normalized, because $y$ is odd; our engine throws on it. And it gives $\alpha = 12.606$, $\bar\alpha = 5.394$, totally positive, so $d = 1$. Same pair, opposite bit.}

\Kle{You found a counterexample to your own source's worked example and you shipped it in the app rather than quietly fixing the example.}

\Stu{The invariant is fine. The definition needs the normalization clause and mere primitivity is not enough. That was worth more than the example was.}

\Kle{It is the single most creditable thing on this tab and I will say so plainly. Now let me break it. ``Every one of the 10 normalized witnesses with $z \le 12$.'' What is the bound on $y$?}

\Stu{The scan runs $y$ from $-120$ to $120$.}

\Kle{Which is not stated. You have declared one of two bounds and phrased the result as ``every one.'' That is the shape of an exhaustive claim over a range you have only half-described. It does not matter that the conclusion is almost certainly right --- on a tab whose entire subject is labelling what is checked and what is not, an undeclared search bound is exactly the wrong omission.}

\Stu{Conceded without qualification. The prose should carry both bounds or neither.}

\begin{knote}
$d$-bit box: three blocks, three statuses, labelled --- empirical (8091 pairs, quoted not computed), computed live (13, 61 totally positive, $d = 1$), and the caution. The caution is the best work in the app: they checked their own source's illustrative witness $(9,1,2)$, found it unnormalized ($y$ odd), found it flips the bit --- same pair, opposite bit --- and published it. That is a laboratory, not a demo.

Against it: search bound half-declared ($z \le 12$ stated, $|y| \le 120$ not). ``Every one of the 10'' is an exhaustiveness claim over an undescribed box. And $d \in \{1,3\}$ is never explained; ``a Stark rank'' is three words for a theory.
\end{knote}

\walkfig{x6-cubictable}{Ten rows: what changes at $\ell = 3$.}

\Stu{Second thread. Everything else in the app is mod 2 --- values in $\{\pm 1\}$, double covers, deck groups $D_4$ and $N_4(\mathbb{F}_2)$. The framework is not restricted to $\ell = 2$. This table is the audit: ten rows, what the $\ell = 2$ column says and what the $\ell = 3$ column says. Value group goes from $\{\pm 1\} = \mu_2$ to $\mu_3 = \{1, \omega, \omega^2\}$; base field from $\mathbb{Q}$ to $k = \mathbb{Q}(\zeta_3) = \mathbb{Q}(\sqrt{-3})$; orientability convention from $p \equiv 1 \pmod 4$ to $p \equiv 1 \pmod 3$, which is what makes $p$ split in $k$; normalization from Redei's $x - y \equiv 1 \pmod 4$ to \emph{primary}, $\pi \equiv 2 \pmod 3$ in $\mathbb{Z}[\omega]$; double cover to 3-fold cyclic cover; quadratic reciprocity to cubic; $N_3(\mathbb{F}_2) \cong D_4$ of order 8 to $N_3(\mathbb{F}_3)$ of order 27 with centre $\mathbb{Z}/3$.}

\Kle{Row eight. ``Chebotarev share per class at the centre --- one half each, one third each.'' Your Q and A tab told me the Ladder's Chebotarev densities were one eighth. Which is it?}

\Stu{Both, and the row is compressed to the point of being misleading. On the Ladder the density of the identity in $D_4$ is $1/8$ and the density of the central corner matrix is $1/8$ --- densities in the whole group. This row is conditional: given that Frobenius has landed in the centre, which for $D_4$ is $\mathbb{Z}/2$ of order two, each of the two elements takes half. For $N_3(\mathbb{F}_3)$ the centre is $\mathbb{Z}/3$ and each takes a third. Same theorem, different denominator.}

\Kle{And the word ``conditional'' is not in the row.}

\Stu{No.}

\Kle{Then a reader who has just come from the Ladder tab sees $1/8$ turn into $1/2$ and has to reconstruct your denominator. Two extra words fixes it. Second point --- the last row calls it ``the $d$-defect,'' and the box we just read calls it ``a governing bit $d\{p,q\}$,'' and the crack above that calls the same territory ``depth-four.'' Three vocabularies for one thread on one tab.}

\Stu{Fair. The table row was written to be scannable and it drifted.}

\Kle{The last two rows are in red and say ``none here'' and ``no real place.'' That is the right instinct --- the failures are in the table rather than in a footnote, so I cannot read the comparison without reading the failures. Keep that.}

\walkfig{x7-mod3topo}{The topology mod 3: three dial positions, $|{\rm deck}| = 3 = $ loops times winding.}

\Stu{Nothing changes in kind, only the modulus. In the 3-fold cyclic cover of $S^3 \setminus K_1$, a loop with monodromy of order $d$ has $3/d$ preimage components each winding $d$ times, so the dial is $\mathrm{lk} \bmod 3$ with three positions instead of two. $\mathrm{lk} \equiv 0$: order 1, three loops winding 1, orbit factor $(1-t)^{-3}$. Otherwise: order 3, one loop winding 3, factor $(1-t^3)^{-1}$. And $3 = \text{loops} \times \text{winding}$ in every row, computed.}

\Kle{This is the row of the dictionary I have no complaint about, because it is just covering space theory and covering space theory does not care what $\ell$ is. Index times sheets. It would be strange if this broke.}

\Stu{Which is why it is the part I am most confident of and least impressed by.}

\Kle{Correct posture. Small thing: you have five rows for three positions. Rows for $\mathrm{lk} = 3$ and $\mathrm{lk} = 4$ are literally rows 0 and 1 again. I assume that is deliberate, to show periodicity --- but nothing on the table says so, so it looks at first glance like the dial has five positions, three sentences after you told me it has three.}

\Stu{It is a periodicity demonstration with no label saying it is one.}

\Kle{And the note at the bottom --- $\bar\mu(123) \bmod 3$ landing at the corner of $N_3(\mathbb{F}_3)$, ``described not computed,'' no mod-3 Magnus pipeline, so no number asserted. Good. That is the sentence I would have demanded if it were not there.}

\walkfig{x8-cubicworked}{The Eisenstein arithmetic, computed live: primary associates and the cubic residue symbol both ways.}

\Stu{Live, at load. Both 7 and 61 are $\equiv 1 \pmod 3$, so both split in $\mathbb{Z}[\omega]$. The app finds a prime above each --- $-3 - 2\omega$ above 7, $-9 - 5\omega$ above 61 --- and then normalizes each to its primary associate by multiplying through the six units until $\pi \equiv 2$ and the $\omega$-coefficient $\equiv 0 \bmod 3$. That gives $\pi_7 = -1 - 3\omega$ and $\pi_{61} = -4 - 9\omega$, and the norms are recomputed: 7 and 61. Then the symbol $\chi_\pi(\alpha) \equiv \alpha^{(N(\pi)-1)/3} \pmod \pi$, valued in $\mu_3$, computed in the residue field, with $\omega \mapsto 2 \bmod 7$ --- and 2 does satisfy $x^2 + x + 1 \equiv 0 \bmod 7$, which the app checks. Result: $\chi_{\pi_7}(\pi_{61}) = 1$ and $\chi_{\pi_{61}}(\pi_7) = 1$. Equal. Cubic reciprocity.}

\Kle{Both of them are 1.}

\Stu{Yes.}

\Kle{So your worked instance of cubic reciprocity is the identity equals the identity. You have verified $1 = 1$ and put a green tick on it. If your implementation of $\chi$ returned the constant 1 for every input, this card would look exactly the same.}

\Stu{That is right and it is a real weakness of the worked example. The evidence that the implementation does anything lives one paragraph down, in the census, where the values are not all trivial.}

\Kle{Then swap them, or add a second row to this card with a nontrivial pair. Choosing a degenerate worked example is a strange thing to do on the one tab where you have made a virtue of finding your own degenerate cases.}

\Stu{The one thing I would keep is the gloss under it: both trivial means 7 and 61 are cubically unlinked --- 61 is a cube mod 7 and 7 is a cube mod 61 --- the mod-3 analogue of Legendre symbol $+1$. That sentence is what makes the number mean something rather than just be equal.}

\Kle{Keep the sentence, change the primes. And the closing line is right: symmetry is what earns the name linking number. That is a real principle and it is the reason the arithmetic side is allowed to say ``linking'' at all --- $\mathrm{lk}(K_1, K_2) = \mathrm{lk}(K_2, K_1)$ is not a convention, it is a theorem, and so is this.}

\walkfig{x9-census}{The 171-pair cubic reciprocity census, computed at load.}

\Stu{All 19 primes $\equiv 1 \pmod 3$ below 200, primary-normalized, and all 171 unordered pairs. Reciprocity $\chi_\pi(\theta) = \chi_\theta(\pi)$ holds 171 out of 171. Value distribution $1 : \omega : \omega^2 = 61 : 57 : 53$, which is 35.7 / 33.3 / 31.0 percent, against the Chebotarev prediction of a third each.}

\Kle{The 171 out of 171 is not evidence of anything mathematical. Eisenstein proved cubic reciprocity. What you have tested is your implementation.}

\Stu{Correct, and the card does say Eisenstein is cited. But it does not say that the 171/171 is a unit test rather than a result, and it should.}

\Kle{Now the distribution, which is the part that is actually a measurement. 61, 57, 53 out of 171. Expected 57 each. What is your standard deviation?}

\Stu{For $n = 171$ at probability one third, the variance is $171 \times (1/3)(2/3) = 38$, so about 6.2.}

\Kle{So your deviations are plus four, zero, minus four --- about two thirds of one standard deviation. Which is to say: perfectly consistent with a third each, and also perfectly consistent with a good many other things, because 171 pairs is nothing. And you have printed 35.7 and 31.0 to one decimal place with no error bar, which invites a reader to see a trend in noise. Either print the standard deviation next to it, or do not print the third significant figure.}

\Stu{That is right. The percentages are more precise than the sample can support.}

\Kle{The closing sentence is honest, at least: your laboratory's real census runs 120 thousand triple symbols and this app cannot compute a single one, because there is no mod-3 triple-symbol construction here. Say that number louder --- the gap between 171 pairs and 120 thousand triples is the whole distance between this tab and the research programme behind it.}

\begin{knote}
Cubic side: primary normalization and $\chi_\pi$ are genuinely computed --- $\pi_7 = -1 - 3\omega$, $\pi_{61} = -4 - 9\omega$, norms recomputed, $\omega \mapsto 2 \bmod 7$ verified against $x^2+x+1$. Solid.

But the worked reciprocity instance is $1 = 1$. Degenerate. A constant function would pass it. The nontrivial evidence is in the census one paragraph below and he knows it.

Census: 171/171 reciprocity tests the code, not the mathematics (Eisenstein proved it) --- not said. Distribution 61:57:53, expected 57, sd $\approx 6.2$; deviations $\approx 0.65$ sd. Consistent, and meaningless at $n = 171$. Printed to one decimal with no error bar.

Covering-space row (mod-3 topology) is unimpeachable and he said so himself --- index times sheets does not care about $\ell$.
\end{knote}

\walkfig{x10-gaps}{The $\ell = 3$ gap list: no Redei construction, no real place, the honesty ladder.}

\Stu{Three gaps. One: no Redei construction. Redei's $\alpha$ is quadratic by design --- it comes from the conic $x^2 - p_1y^2 - p_2z^2 = 0$. The $\ell = 3$ analogue needs a degree-3 Kummer relative-norm solver, which this app does not have, so the mod-3 triple symbol is described here and never computed. Two: the defect has nowhere at infinity to live. At $\ell = 2$ the governing bit is a sign --- a real place, totally real against totally imaginary. At $\ell = 3$ the base field is $\mathbb{Q}(\sqrt{-3})$, imaginary quadratic, zero real embeddings. There is no real place for a sign to sit in, so the conjecture is that the defect migrates to the ramified prime $\lambda = (\sqrt{-3})$ above 3 --- a $\lambda$-adic $d_\lambda$ replacing an archimedean sign. Labelled conjectural, and it is the crux of the programme.}

\Kle{Stop. Gap two says ``zero real embeddings'' and then, in parentheses, ``computed at right.'' Show me where at right that is computed.}

\walkbeat{The student scrolls up to the right-hand column and stops.}

\Stu{It is not. The right-hand column computes primary associates, cubic residue symbols and reciprocity. It never computes a signature. $\mathbb{Q}(\sqrt{-3})$ having no real embeddings is stated in the table row above, not computed anywhere.}

\Kle{So on the tab whose entire discipline is the four-way label --- theorem, computed live, empirical, conjectural --- you have a parenthesis reading ``computed at right'' pointing at a computation that does not exist. That is the one error I would call serious, because it is not a gap in the mathematics, it is a false label on the honesty system itself. Everything else on this tab I can check by reading. That one I could only catch by scrolling.}

\Stu{I have no defence and I am not going to construct one. It is a wrong word in the one place the app cannot afford a wrong word.}

\Kle{Which is, in its way, a compliment to the rest of it --- I had to look at all eleven boxes to find one. Read me the third gap.}

\Stu{The honesty ladder. $\ell = 2$ is proven: Legendre, Redei, Amano, and the one machine-checked depth-four law. The $\ell = 3$ migration to $d_\lambda$ is conjectural. The $p$-adic founding conjecture behind both is speculative. This tab exists to keep those three tiers visibly apart.}

\Kle{Three tiers, and the sentence naming them is set in small grey centred italics at the bottom of a wall of small grey centred italics. I said this at the top of the tab and I will say it once more because it is the same defect three times: your gap lists are your best writing and your worst typography. A reader skims grey italic. You have put the load in the place the eye is trained to skip.}

\begin{knote}
Serious: gap (ii) says $\mathbb{Q}(\sqrt{-3})$ has zero real embeddings ``(computed at right)'' --- nothing at right computes a signature. A false ``computed'' label on the tab whose whole subject is what is computed. Not a mathematical error; a labelling error, which here is worse.

Otherwise the gap list is the right list: no degree-3 Kummer solver so no mod-3 triple symbol, ever, in this app --- said outright. The $d_\lambda$ migration conjecture is clearly marked. The three-tier ladder (proven / conjectural / speculative) is exactly the structure this material needs.

Typography, third time: the three most important paragraphs on the tab are the three least readable. Systematic, not accidental --- shared style class.
\end{knote}

\subsection*{The self-test}

\walkbeat{The student scrolls to the footer and presses Run self-tests. The panel fills in a single frame.}

\walkfig{s1-tests}{The suite, run live: 114 of 114 passing.}

\Stu{114 of 114, and the button relabels itself with the split: 62 computational, 52 rendered-text pins. Two species, and the count says which is which because you cannot judge the suite without knowing the mixture.}

\Kle{Define the second species.}

\Stu{A rendered-text pin reads the actual DOM after render and asserts that the prose says what it is supposed to say. The Q and A one, for instance, asserts there are exactly 14 question items, that all 14 are closed at load, that at least ten jump buttons exist, and that the words ``Mazur,'' ``Artin--Verdier,'' ``Chebotarev'' and ``It is not a functor'' are present in the rendered text.}

\Kle{Then be precise about what that species can and cannot catch. It can catch the prose drifting away from the code. It cannot catch the prose and the code being wrong together, and it cannot catch a theorem being misattributed --- if you cited Milnor 1962 for something Turaev proved, all 52 would stay green.}

\Stu{That is exactly right and it is the limit of the whole suite. The pins catch drift. They do not catch error. What the 62 computational ones catch is different: those assert against expected values that do not come from the same run.}

\Kle{Give me one I can check sitting here.}

\Stu{Line six. \texttt{tonelliShanks(13,937) === [386,551]}.}

\walkbeat{Klein does the arithmetic on the back of the handout.}

\Kle{$386^2 = 148996$. And $937 \times 159 = 148983$. Difference 13. So $386^2 \equiv 13 \pmod{937}$, and the other root has to be $937 - 386 = 551$. Your square root is a square root.}

\Stu{And that is the point of that species: you did not have to trust the implementation, you had to trust multiplication.}

\Kle{Which is the correct place to put the trust. Now --- what happens if I break one?}

\Stu{The failing line turns red and bold, and the button text changes to the passing count out of 114 followed by ``see failures.'' It does not tell you which line, so you scroll.}

\Kle{Through 114 monospace lines in no particular order, with no grouping by tab. On a green run that does not matter. On a red run it is the only moment the suite is useful and it is the moment it helps least.}

\Stu{Fair. There is no grouping and no jump-to-failure.}

\walkbeat{Klein runs his finger down the list and stops partway.}

\Kle{``Klein F1 pure.'' ``Klein A2 pure.'' ``Klein A4 pure.'' ``DOM Klein A1/A2.'' ``Klein C8 pure.''}

\Stu{Those are from your earlier review rounds. Each objection you made that was specific enough to be falsifiable got turned into a test with your initials on it, so it cannot silently come back.}

\Kle{A4 is the Chebotarev tally.}

\Stu{427 unramified odd primes below 3000. Identity 50, corner 47, and the three character classes 109, 111, 110 --- the observed frequencies hugging one eighth, one eighth, one quarter, one quarter, one quarter. That is the finite shadow of the density statement the Q and A tab makes in prose, pinned as a number.}

\Kle{And the certificate one, four from the bottom.}

\Stu{$[13,61,269] = -1$ via the certificate $(23,6,1)$, and then the $k=2$ scaling $(46,12,2)$ --- which satisfies the same identity --- would report $+1$, because scaling by $k$ multiplies the value by $(k/q)$ and $(2/269) = -1$. The engine rejects it. That test exists because the symbol is only well defined on normalized certificates, and the failure mode is silent otherwise.}

\Kle{So the suite's most interesting entries are all of the same kind: not ``the code gets the right answer,'' but ``the code refuses the wrong question.'' The rejection tests are worth more than the computation tests, because a wrong number is visible and a silently accepted invalid input is not.}

\Stu{The rejection taxonomy in \texttt{redeiSymbol} was the thing that took longest to get right. Primality, definedness mod 4, the identity, primitivity, $y$ even, $x - y \equiv 1 \pmod 4$, and then the case where $q$ divides a conjugate --- each with its own message, in that order, so a bad input hears the real reason first.}

\Kle{Then here is my summary of the suite, and it is a mixed one. 114 green is not a proof of anything and you have not claimed it is --- the split in the button is you telling me so before I ask, which I note. What it is proof of is that this app is a piece of engineering with a regression net under it, which is not the normal condition of a teaching demo, and which is why I have been able to spend two hours arguing about mathematics instead of about whether the numbers are stable.}

\begin{knote}
Ran the suite live: 114/114, split shown on the button as 62 computational, 52 rendered-text pins. Naming the split unprompted is the right instinct --- I did not have to ask what fraction were tautologies.

Spot-checked \texttt{tonelliShanks(13,937) = [386,551]} by hand: $386^2 = 148996 = 937 \cdot 159 + 13$, and $551 = 937 - 386$. Passes.

Limits, which he stated before I did: the 52 pins catch prose/code drift, not error --- a misattributed theorem stays green. The 62 computational assert against constants chosen off-run, so they are checkable, and I checked one.

Best entries are the rejection tests, not the computation tests: the scaled certificate $(46,12,2)$ satisfying the same identity and reporting $+1$ via $(2/269) = -1$, refused. A wrong number is visible; a silently accepted invalid input is not.

Weakness: 114 flat monospace lines, no grouping, no jump-to-failure. Useless exactly when it matters.

My initials are in five test names. That is either flattery or good engineering practice and I suspect it is both.
\end{knote}\clearpage
\section*{Post-mortem}
\addcontentsline{toc}{section}{Post-mortem}

\walkbeat{The laptop is closed. Klein has his notebook open on the table and is going back through it with a pen, not turning pages in order.}

\Kle{Sit down. I want to do this while it is fresh, and I want to separate three things that people always run together: is it correct, is it honest, and is it legible. You did very well on the first, better than most published expositions on the second, and you have a systematic problem on the third.}

\Stu{I would rather have it in that order than the reverse.}

\Kle{Correctness first, because it is the shortest conversation. I attacked the arithmetic in four places and it held every time. The Euler product for $\zeta_{k_{13,61}}$ agreeing with the ideal-sum face to seven decimals is not a coincidence you can fake --- those are genuinely different computations over the same data. The order rule brute-forced over all sixty-four elements of $N_4(\mathbb{F}_2)$ is the right way to present a lemma that nobody wants to check by hand. And $101$ being irregular is a fact about the world that your app found sitting inside its own cast list. That is the best moment in the presentation and you did not oversell it, which is why it landed.}

\Stu{The $101$ was not planted. I want that on the record, because it is the one moment where the dictionary generated a question rather than illustrating an answer.}

\Kle{It is on the record. Now honesty, where you are mostly excellent and where your one real failure lives. Your discipline is that every number is computed at load and everything else carries a name and a year, and that \emph{computed}, \emph{cited}, and \emph{analogy} never trade places. You enforce that almost everywhere --- the Milnor-model gap on $B_4$, the withheld gloss at $(3,7)$, the blank arithmetic cell in the five-component row, the refusal to assert a $449$ factor when the conjugate check fails. Those are not decorations. They are the reason I believed the rest.}

\Stu{And the failure.}

\Kle{On the Experimental tab, in the gap list, you write that $\mathbb{Q}(\sqrt{-3})$ has zero real embeddings and you append the words \emph{computed at right}. Nothing at right computes a signature. That column computes primary associates, cubic residue symbols, and a reciprocity check. The claim is true --- it is imaginary quadratic, a first-year fact --- but the label is false, and it is false on the one tab whose entire subject is the difference between what you checked and what you are repeating. That is the single thing I would fix before you show this to anybody else.}

\Stu{It is a false provenance label, not a false statement, which somehow makes it worse.}

\Kle{Considerably worse, because your test suite cannot catch it. Which brings me to the general point about those one hundred fourteen checks. Fifty-two of them pin rendered prose against the page. That catches drift --- code changing while a sentence stays behind. It cannot catch code and prose being confidently wrong together, and it cannot catch a misattributed citation at all. Your suite would stay green through a wrong year, a wrong author, or a mislabelled provenance. Say that somewhere, because right now the green number invites more trust than it can carry.}

\Stu{The button says sixty-two computational and fifty-two rendered-text pins, but it does not say what the second number cannot do.}

\Kle{Add the sentence. Now legibility, and this is the systematic one. You have a recurring habit: the most load-bearing justification on a screen is the least visible thing on it. Artin-Verdier is the wall your entire premise leans on, and it is inside a hover tooltip --- an unaided reader sees ``behaves three-dimensionally'' followed by a parenthesis. The point that $\gamma$ is analogous to $\varphi_*$ and \emph{not} to Frobenius --- which your own tooltip calls the commonest error in reading the dictionary --- is also in a tooltip. And on the Experimental tab, your shared \texttt{honest} class renders the three longest and most important paragraphs in small, muted, centred italics. The class is named for its epistemic function and styled as a disclaimer.}

\Stu{The styling was meant to signal ``this is a caveat, read carefully.''}

\Kle{It signals ``this is fine print, skip.'' Those are opposite instructions and the typography wins. And it interacts badly with the tab banner: you have quarantined behind ``nothing here is load-bearing'' what I think is the sharpest structural idea in the whole app --- that feeding the Lefschetz formula a ramified generator gives you knot invariants and feeding it Frobenius gives you link invariants and Artin $L$-functions. That is not frontier speculation. That is a clean observation about the two halves of your own app, and it explains what the other seven tabs are doing. The empirical $d$-bit law belongs behind that banner. The pivot does not.}

\Stu{So the banner is measuring the wrong thing. It marks the tab, when it should mark the claims.}

\Kle{Precisely. Mark the claims; you already have a vocabulary for it. Three more, quickly. Your terminology drifts inside single screens --- the same thread is a \emph{$d$-defect}, a \emph{governing bit}, and \emph{the machine-checked depth-four case}, and \emph{depth-four} is never defined while \emph{depth two} is defined ten centimetres away. Second, the trefoil table computes orders for the \emph{branched} cyclic covers and the word branched never appears; a reader who takes ``cyclic covers'' literally computes the wrong object. Third, and this one is a gift rather than a complaint: those orders $3, 4, 3, 1, \infty$ are the Brieskorn spheres $\Sigma(2,3,n)$. The $n = 5$ entry is the Poincare homology sphere. You print it as the integer $1$ and say nothing. Every topologist you show this to will recognise the sequence and wonder why you did not.}

\Stu{That is the payoff I did not know I had.}

\Kle{You have several. The other one is on Pairs: choosing both primes congruent to $3$ mod $4$ makes the $2$-adic Hilbert factor appear as a \emph{term} in the product formula instead of an exception to it --- which is exactly the argument that the $1$ mod $4$ convention is a choice to work where one local term vanishes, not a technicality. Every piece of that is on your screen. It is never assembled into one sentence, so the reader has to build it from the warning line, the product note, and a remark two tabs back.}

\Stu{What would you have me do first?}

\Kle{Fix the false label. Then decide, once, where load-bearing prose lives --- and I would say: never in a tooltip, never in the \texttt{honest} class, never centred and muted. Then unify the vocabulary. The rest is polish, and the polish is good enough that I noticed it. Two things I will say plainly because you have absorbed an hour of criticism without flinching: the surface ladder is the best piece of mathematical exposition in the app, and the three-step rung-two figure is what made the intersection curve finally sit still for me. And your refusals --- the blank cell, the withheld gloss, the model-level caveat --- are what convinced me the numbers were worth checking in the first place.}

\Stu{And the thing you would still refuse to sign.}

\Kle{The depth-three proposal, obviously, which you already label a proposal. And I would not yet say the dictionary explains anything on the topological side. It organises, it predicts where to look, and once --- with $101$ --- it produced a question. That is a real result for an exhibit. It is not a theorem, and your app is careful never to say it is.}

\begin{knote}
Final. Correct throughout; four attacks, four holds. Honest by construction, one false provenance label (Exp/gap list, ``computed at right''). Systematic legibility fault: load-bearing text lives in tooltips and in the honest class, i.e. the app whispers exactly where it should speak. Vocabulary drifts. Missed payoffs: Brieskorn, and the 2-adic-factor-as-term argument. Best things: the surface ladder, the rung-2 storyboard, the refusals, and $101$. Would I show this to a graduate class? Yes --- after the label is fixed and the tooltips are promoted.
\end{knote}

\subsection*{The punch list}

Ordered by what Klein would fix first. Every item traces to a specific moment in the walkthrough above.

\textbf{P1 --- honesty defects (fix before the next showing).}
\begin{enumerate}[leftmargin=1.4em,itemsep=2pt]
\item \textbf{False provenance label.} The $\ell = 3$ gap list claims $\mathbb{Q}(\sqrt{-3})$ has zero real embeddings ``computed at right''; nothing computes a signature. Either compute it or restate it as the standard fact it is.
\item \textbf{The suite oversells itself.} Add one sentence to the self-test panel: the $52$ rendered-text pins detect drift between code and prose, not code and prose being wrong together, and detect nothing about citations.
\item \textbf{Undisclosed search bounds.} The $(13,17)$ witness scan says ``every one of the $10$ normalized witnesses with $z \le 12$'' but also silently bounds $|y| \le 120$. State the whole box.
\item \textbf{Typographic equality between computed and asserted rows.} In the Hilbert table, the ``all other places are $+1$'' row is a cited theorem about an infinite family and is set exactly like the two enumerated rows beside it.
\item \textbf{A degenerate check presented as evidence.} The worked cubic instance is $1 = 1$ with a green tick; a constant-$1$ implementation would render identically. Pick a pair with a nontrivial symbol, and keep the census.
\end{enumerate}

\textbf{P2 --- load-bearing prose is hidden (the systematic fault).}
\begin{enumerate}[leftmargin=1.4em,itemsep=2pt,resume]
\item \textbf{Promote Artin-Verdier} out of its tooltip into the visible why-box: the trace isomorphism, $\mathbb{G}_m$ as orientation sheaf, and the real-place correction are the justification for the app's premise.
\item \textbf{Promote the generator distinction} --- $\gamma \sim \varphi_*$, not $\gamma \sim \mathrm{Frob}$ --- into body text. The app's own tooltip calls conflating them the commonest error in reading the dictionary.
\item \textbf{Restyle the \texttt{honest} class.} Small, muted, centred italics reads as fine print; these are the app's most careful paragraphs. Same treatment, opposite signal.
\item \textbf{Re-scope the Experimental banner.} Label the claims, not the tab. The ramified/unramified pivot is an observation about the existing seven tabs and should not sit behind ``nothing here is load-bearing''.
\end{enumerate}

\textbf{P3 --- vocabulary and definitions.}
\begin{enumerate}[leftmargin=1.4em,itemsep=2pt,resume]
\item \textbf{One name per object.} ``$d$-defect'', ``governing bit $d\{p,q\}$'', and ``the machine-checked depth-four case'' are one thread under three names on one screen.
\item \textbf{Define \emph{depth-four}} where it is used, given that \emph{depth two} is defined in the adjacent column; the natural misreading is ``four components''.
\item \textbf{Say \emph{branched}.} The trefoil table computes $\#H_1$ of the branched cyclic covers; $M_n$ is never defined and ``branched'' never appears.
\item \textbf{Define $f_S$ on the Q and A tab}, where it is used twice and defined nowhere.
\end{enumerate}

\textbf{P4 --- state, navigation, and figures.}
\begin{enumerate}[leftmargin=1.4em,itemsep=2pt,resume]
\item \textbf{Stale canvas at $(3,7)$.} A deliberate refusal to sync renders identically to a forgotten update; say which it is.
\item \textbf{The Experimental tab is a navigational dead end} --- zero jump buttons back into the tabs whose claims it reframes; the Q and A tab returns the favour by never linking it.
\item \textbf{Q and A ergonomics.} Fourteen questions, unnumbered, no expand-all, fourteen clicks to read.
\item \textbf{Figure defects.} The arc figure paints two labels under the blue curve; several $\tilde{X}$ labels render with a displaced combining accent; one certificate is joined with an ASCII hyphen beside a typographic minus.
\item \textbf{Column balance} on the Experimental generators row: the right column ends halfway, leaving the tab's most delicate prose beside white space.
\end{enumerate}

\textbf{P5 --- payoffs currently forfeited.}
\begin{enumerate}[leftmargin=1.4em,itemsep=2pt,resume]
\item \textbf{Name the Brieskorn spheres.} $3, 4, 3, 1, \infty$ are $\Sigma(2,3,n)$; the $1$ is the Poincare homology sphere. Every topologist will recognise the sequence.
\item \textbf{Assemble the $2$-adic argument} into one sentence on Pairs: both primes $\equiv 3 \pmod 4$ makes the $2$-adic factor a term rather than an exception, which is the real content of the $1$ mod $4$ convention.
\item \textbf{Flag the degenerate agreement.} The knot-column match is exhibited only where $\lambda = \mu = \nu = 0$; agreement by mutual vanishing deserves the caveat in the display, not just in the surrounding prose.
\item \textbf{Give the census an error bar.} At $n = 171$, $61:57:53$ sits about $0.65$ standard deviations from uniform; printing one decimal place with no interval claims more than the sample supports, and the self-test pins those three integers as exact constants.
\end{enumerate}

\walkbeat{Klein closes the notebook and pushes it across the table.}

\Kle{Take the notes. I want the label fixed and the tooltips promoted; the rest you can argue with me about. And put the Brieskorn line in --- that one costs you nothing and buys you the room.}

\end{document}
